\documentclass[11pt]{article}

    \usepackage[breakable]{tcolorbox}
    \usepackage{parskip} % Stop auto-indenting (to mimic markdown behaviour)
    

    % Basic figure setup, for now with no caption control since it's done
    % automatically by Pandoc (which extracts ![](path) syntax from Markdown).
    \usepackage{graphicx}
    % Maintain compatibility with old templates. Remove in nbconvert 6.0
    \let\Oldincludegraphics\includegraphics
    % Ensure that by default, figures have no caption (until we provide a
    % proper Figure object with a Caption API and a way to capture that
    % in the conversion process - todo).
    \usepackage{caption}
    \DeclareCaptionFormat{nocaption}{}
    \captionsetup{format=nocaption,aboveskip=0pt,belowskip=0pt}

    \usepackage{float}
    \floatplacement{figure}{H} % forces figures to be placed at the correct location
    \usepackage{xcolor} % Allow colors to be defined
    \usepackage{enumerate} % Needed for markdown enumerations to work
    \usepackage{geometry} % Used to adjust the document margins
    \usepackage{amsmath} % Equations
    \usepackage{amssymb} % Equations
    \usepackage{textcomp} % defines textquotesingle
    % Hack from http://tex.stackexchange.com/a/47451/13684:
    \AtBeginDocument{%
        \def\PYZsq{\textquotesingle}% Upright quotes in Pygmentized code
    }
    \usepackage{upquote} % Upright quotes for verbatim code
    \usepackage{eurosym} % defines \euro

    \usepackage{iftex}
    \ifPDFTeX
        \usepackage[T1]{fontenc}
        \IfFileExists{alphabeta.sty}{
              \usepackage{alphabeta}
          }{
              \usepackage[mathletters]{ucs}
              \usepackage[utf8x]{inputenc}
          }
    \else
        \usepackage{fontspec}
        \usepackage{unicode-math}
    \fi

    \usepackage{fancyvrb} % verbatim replacement that allows latex
    \usepackage{grffile} % extends the file name processing of package graphics 
                         % to support a larger range
    \makeatletter % fix for old versions of grffile with XeLaTeX
    \@ifpackagelater{grffile}{2019/11/01}
    {
      % Do nothing on new versions
    }
    {
      \def\Gread@@xetex#1{%
        \IfFileExists{"\Gin@base".bb}%
        {\Gread@eps{\Gin@base.bb}}%
        {\Gread@@xetex@aux#1}%
      }
    }
    \makeatother
    \usepackage[Export]{adjustbox} % Used to constrain images to a maximum size
    \adjustboxset{max size={0.9\linewidth}{0.9\paperheight}}

    % The hyperref package gives us a pdf with properly built
    % internal navigation ('pdf bookmarks' for the table of contents,
    % internal cross-reference links, web links for URLs, etc.)
    \usepackage{hyperref}
    % The default LaTeX title has an obnoxious amount of whitespace. By default,
    % titling removes some of it. It also provides customization options.
    \usepackage{titling}
    \usepackage{longtable} % longtable support required by pandoc >1.10
    \usepackage{booktabs}  % table support for pandoc > 1.12.2
    \usepackage{array}     % table support for pandoc >= 2.11.3
    \usepackage{calc}      % table minipage width calculation for pandoc >= 2.11.1
    \usepackage[inline]{enumitem} % IRkernel/repr support (it uses the enumerate* environment)
    \usepackage[normalem]{ulem} % ulem is needed to support strikethroughs (\sout)
                                % normalem makes italics be italics, not underlines
    \usepackage{mathrsfs}
    

    
    % Colors for the hyperref package
    \definecolor{urlcolor}{rgb}{0,.145,.698}
    \definecolor{linkcolor}{rgb}{.71,0.21,0.01}
    \definecolor{citecolor}{rgb}{.12,.54,.11}

    % ANSI colors
    \definecolor{ansi-black}{HTML}{3E424D}
    \definecolor{ansi-black-intense}{HTML}{282C36}
    \definecolor{ansi-red}{HTML}{E75C58}
    \definecolor{ansi-red-intense}{HTML}{B22B31}
    \definecolor{ansi-green}{HTML}{00A250}
    \definecolor{ansi-green-intense}{HTML}{007427}
    \definecolor{ansi-yellow}{HTML}{DDB62B}
    \definecolor{ansi-yellow-intense}{HTML}{B27D12}
    \definecolor{ansi-blue}{HTML}{208FFB}
    \definecolor{ansi-blue-intense}{HTML}{0065CA}
    \definecolor{ansi-magenta}{HTML}{D160C4}
    \definecolor{ansi-magenta-intense}{HTML}{A03196}
    \definecolor{ansi-cyan}{HTML}{60C6C8}
    \definecolor{ansi-cyan-intense}{HTML}{258F8F}
    \definecolor{ansi-white}{HTML}{C5C1B4}
    \definecolor{ansi-white-intense}{HTML}{A1A6B2}
    \definecolor{ansi-default-inverse-fg}{HTML}{FFFFFF}
    \definecolor{ansi-default-inverse-bg}{HTML}{000000}

    % common color for the border for error outputs.
    \definecolor{outerrorbackground}{HTML}{FFDFDF}

    % commands and environments needed by pandoc snippets
    % extracted from the output of `pandoc -s`
    \providecommand{\tightlist}{%
      \setlength{\itemsep}{0pt}\setlength{\parskip}{0pt}}
    \DefineVerbatimEnvironment{Highlighting}{Verbatim}{commandchars=\\\{\}}
    % Add ',fontsize=\small' for more characters per line
    \newenvironment{Shaded}{}{}
    \newcommand{\KeywordTok}[1]{\textcolor[rgb]{0.00,0.44,0.13}{\textbf{{#1}}}}
    \newcommand{\DataTypeTok}[1]{\textcolor[rgb]{0.56,0.13,0.00}{{#1}}}
    \newcommand{\DecValTok}[1]{\textcolor[rgb]{0.25,0.63,0.44}{{#1}}}
    \newcommand{\BaseNTok}[1]{\textcolor[rgb]{0.25,0.63,0.44}{{#1}}}
    \newcommand{\FloatTok}[1]{\textcolor[rgb]{0.25,0.63,0.44}{{#1}}}
    \newcommand{\CharTok}[1]{\textcolor[rgb]{0.25,0.44,0.63}{{#1}}}
    \newcommand{\StringTok}[1]{\textcolor[rgb]{0.25,0.44,0.63}{{#1}}}
    \newcommand{\CommentTok}[1]{\textcolor[rgb]{0.38,0.63,0.69}{\textit{{#1}}}}
    \newcommand{\OtherTok}[1]{\textcolor[rgb]{0.00,0.44,0.13}{{#1}}}
    \newcommand{\AlertTok}[1]{\textcolor[rgb]{1.00,0.00,0.00}{\textbf{{#1}}}}
    \newcommand{\FunctionTok}[1]{\textcolor[rgb]{0.02,0.16,0.49}{{#1}}}
    \newcommand{\RegionMarkerTok}[1]{{#1}}
    \newcommand{\ErrorTok}[1]{\textcolor[rgb]{1.00,0.00,0.00}{\textbf{{#1}}}}
    \newcommand{\NormalTok}[1]{{#1}}
    
    % Additional commands for more recent versions of Pandoc
    \newcommand{\ConstantTok}[1]{\textcolor[rgb]{0.53,0.00,0.00}{{#1}}}
    \newcommand{\SpecialCharTok}[1]{\textcolor[rgb]{0.25,0.44,0.63}{{#1}}}
    \newcommand{\VerbatimStringTok}[1]{\textcolor[rgb]{0.25,0.44,0.63}{{#1}}}
    \newcommand{\SpecialStringTok}[1]{\textcolor[rgb]{0.73,0.40,0.53}{{#1}}}
    \newcommand{\ImportTok}[1]{{#1}}
    \newcommand{\DocumentationTok}[1]{\textcolor[rgb]{0.73,0.13,0.13}{\textit{{#1}}}}
    \newcommand{\AnnotationTok}[1]{\textcolor[rgb]{0.38,0.63,0.69}{\textbf{\textit{{#1}}}}}
    \newcommand{\CommentVarTok}[1]{\textcolor[rgb]{0.38,0.63,0.69}{\textbf{\textit{{#1}}}}}
    \newcommand{\VariableTok}[1]{\textcolor[rgb]{0.10,0.09,0.49}{{#1}}}
    \newcommand{\ControlFlowTok}[1]{\textcolor[rgb]{0.00,0.44,0.13}{\textbf{{#1}}}}
    \newcommand{\OperatorTok}[1]{\textcolor[rgb]{0.40,0.40,0.40}{{#1}}}
    \newcommand{\BuiltInTok}[1]{{#1}}
    \newcommand{\ExtensionTok}[1]{{#1}}
    \newcommand{\PreprocessorTok}[1]{\textcolor[rgb]{0.74,0.48,0.00}{{#1}}}
    \newcommand{\AttributeTok}[1]{\textcolor[rgb]{0.49,0.56,0.16}{{#1}}}
    \newcommand{\InformationTok}[1]{\textcolor[rgb]{0.38,0.63,0.69}{\textbf{\textit{{#1}}}}}
    \newcommand{\WarningTok}[1]{\textcolor[rgb]{0.38,0.63,0.69}{\textbf{\textit{{#1}}}}}
    
    
    % Define a nice break command that doesn't care if a line doesn't already
    % exist.
    \def\br{\hspace*{\fill} \\* }
    % Math Jax compatibility definitions
    \def\gt{>}
    \def\lt{<}
    \let\Oldtex\TeX
    \let\Oldlatex\LaTeX
    \renewcommand{\TeX}{\textrm{\Oldtex}}
    \renewcommand{\LaTeX}{\textrm{\Oldlatex}}
    % Document parameters
    % Document title
    \title{thesis\_project}
    
    
    
    
    
% Pygments definitions
\makeatletter
\def\PY@reset{\let\PY@it=\relax \let\PY@bf=\relax%
    \let\PY@ul=\relax \let\PY@tc=\relax%
    \let\PY@bc=\relax \let\PY@ff=\relax}
\def\PY@tok#1{\csname PY@tok@#1\endcsname}
\def\PY@toks#1+{\ifx\relax#1\empty\else%
    \PY@tok{#1}\expandafter\PY@toks\fi}
\def\PY@do#1{\PY@bc{\PY@tc{\PY@ul{%
    \PY@it{\PY@bf{\PY@ff{#1}}}}}}}
\def\PY#1#2{\PY@reset\PY@toks#1+\relax+\PY@do{#2}}

\@namedef{PY@tok@w}{\def\PY@tc##1{\textcolor[rgb]{0.73,0.73,0.73}{##1}}}
\@namedef{PY@tok@c}{\let\PY@it=\textit\def\PY@tc##1{\textcolor[rgb]{0.24,0.48,0.48}{##1}}}
\@namedef{PY@tok@cp}{\def\PY@tc##1{\textcolor[rgb]{0.61,0.40,0.00}{##1}}}
\@namedef{PY@tok@k}{\let\PY@bf=\textbf\def\PY@tc##1{\textcolor[rgb]{0.00,0.50,0.00}{##1}}}
\@namedef{PY@tok@kp}{\def\PY@tc##1{\textcolor[rgb]{0.00,0.50,0.00}{##1}}}
\@namedef{PY@tok@kt}{\def\PY@tc##1{\textcolor[rgb]{0.69,0.00,0.25}{##1}}}
\@namedef{PY@tok@o}{\def\PY@tc##1{\textcolor[rgb]{0.40,0.40,0.40}{##1}}}
\@namedef{PY@tok@ow}{\let\PY@bf=\textbf\def\PY@tc##1{\textcolor[rgb]{0.67,0.13,1.00}{##1}}}
\@namedef{PY@tok@nb}{\def\PY@tc##1{\textcolor[rgb]{0.00,0.50,0.00}{##1}}}
\@namedef{PY@tok@nf}{\def\PY@tc##1{\textcolor[rgb]{0.00,0.00,1.00}{##1}}}
\@namedef{PY@tok@nc}{\let\PY@bf=\textbf\def\PY@tc##1{\textcolor[rgb]{0.00,0.00,1.00}{##1}}}
\@namedef{PY@tok@nn}{\let\PY@bf=\textbf\def\PY@tc##1{\textcolor[rgb]{0.00,0.00,1.00}{##1}}}
\@namedef{PY@tok@ne}{\let\PY@bf=\textbf\def\PY@tc##1{\textcolor[rgb]{0.80,0.25,0.22}{##1}}}
\@namedef{PY@tok@nv}{\def\PY@tc##1{\textcolor[rgb]{0.10,0.09,0.49}{##1}}}
\@namedef{PY@tok@no}{\def\PY@tc##1{\textcolor[rgb]{0.53,0.00,0.00}{##1}}}
\@namedef{PY@tok@nl}{\def\PY@tc##1{\textcolor[rgb]{0.46,0.46,0.00}{##1}}}
\@namedef{PY@tok@ni}{\let\PY@bf=\textbf\def\PY@tc##1{\textcolor[rgb]{0.44,0.44,0.44}{##1}}}
\@namedef{PY@tok@na}{\def\PY@tc##1{\textcolor[rgb]{0.41,0.47,0.13}{##1}}}
\@namedef{PY@tok@nt}{\let\PY@bf=\textbf\def\PY@tc##1{\textcolor[rgb]{0.00,0.50,0.00}{##1}}}
\@namedef{PY@tok@nd}{\def\PY@tc##1{\textcolor[rgb]{0.67,0.13,1.00}{##1}}}
\@namedef{PY@tok@s}{\def\PY@tc##1{\textcolor[rgb]{0.73,0.13,0.13}{##1}}}
\@namedef{PY@tok@sd}{\let\PY@it=\textit\def\PY@tc##1{\textcolor[rgb]{0.73,0.13,0.13}{##1}}}
\@namedef{PY@tok@si}{\let\PY@bf=\textbf\def\PY@tc##1{\textcolor[rgb]{0.64,0.35,0.47}{##1}}}
\@namedef{PY@tok@se}{\let\PY@bf=\textbf\def\PY@tc##1{\textcolor[rgb]{0.67,0.36,0.12}{##1}}}
\@namedef{PY@tok@sr}{\def\PY@tc##1{\textcolor[rgb]{0.64,0.35,0.47}{##1}}}
\@namedef{PY@tok@ss}{\def\PY@tc##1{\textcolor[rgb]{0.10,0.09,0.49}{##1}}}
\@namedef{PY@tok@sx}{\def\PY@tc##1{\textcolor[rgb]{0.00,0.50,0.00}{##1}}}
\@namedef{PY@tok@m}{\def\PY@tc##1{\textcolor[rgb]{0.40,0.40,0.40}{##1}}}
\@namedef{PY@tok@gh}{\let\PY@bf=\textbf\def\PY@tc##1{\textcolor[rgb]{0.00,0.00,0.50}{##1}}}
\@namedef{PY@tok@gu}{\let\PY@bf=\textbf\def\PY@tc##1{\textcolor[rgb]{0.50,0.00,0.50}{##1}}}
\@namedef{PY@tok@gd}{\def\PY@tc##1{\textcolor[rgb]{0.63,0.00,0.00}{##1}}}
\@namedef{PY@tok@gi}{\def\PY@tc##1{\textcolor[rgb]{0.00,0.52,0.00}{##1}}}
\@namedef{PY@tok@gr}{\def\PY@tc##1{\textcolor[rgb]{0.89,0.00,0.00}{##1}}}
\@namedef{PY@tok@ge}{\let\PY@it=\textit}
\@namedef{PY@tok@gs}{\let\PY@bf=\textbf}
\@namedef{PY@tok@ges}{\let\PY@bf=\textbf\let\PY@it=\textit}
\@namedef{PY@tok@gp}{\let\PY@bf=\textbf\def\PY@tc##1{\textcolor[rgb]{0.00,0.00,0.50}{##1}}}
\@namedef{PY@tok@go}{\def\PY@tc##1{\textcolor[rgb]{0.44,0.44,0.44}{##1}}}
\@namedef{PY@tok@gt}{\def\PY@tc##1{\textcolor[rgb]{0.00,0.27,0.87}{##1}}}
\@namedef{PY@tok@err}{\def\PY@bc##1{{\setlength{\fboxsep}{\string -\fboxrule}\fcolorbox[rgb]{1.00,0.00,0.00}{1,1,1}{\strut ##1}}}}
\@namedef{PY@tok@kc}{\let\PY@bf=\textbf\def\PY@tc##1{\textcolor[rgb]{0.00,0.50,0.00}{##1}}}
\@namedef{PY@tok@kd}{\let\PY@bf=\textbf\def\PY@tc##1{\textcolor[rgb]{0.00,0.50,0.00}{##1}}}
\@namedef{PY@tok@kn}{\let\PY@bf=\textbf\def\PY@tc##1{\textcolor[rgb]{0.00,0.50,0.00}{##1}}}
\@namedef{PY@tok@kr}{\let\PY@bf=\textbf\def\PY@tc##1{\textcolor[rgb]{0.00,0.50,0.00}{##1}}}
\@namedef{PY@tok@bp}{\def\PY@tc##1{\textcolor[rgb]{0.00,0.50,0.00}{##1}}}
\@namedef{PY@tok@fm}{\def\PY@tc##1{\textcolor[rgb]{0.00,0.00,1.00}{##1}}}
\@namedef{PY@tok@vc}{\def\PY@tc##1{\textcolor[rgb]{0.10,0.09,0.49}{##1}}}
\@namedef{PY@tok@vg}{\def\PY@tc##1{\textcolor[rgb]{0.10,0.09,0.49}{##1}}}
\@namedef{PY@tok@vi}{\def\PY@tc##1{\textcolor[rgb]{0.10,0.09,0.49}{##1}}}
\@namedef{PY@tok@vm}{\def\PY@tc##1{\textcolor[rgb]{0.10,0.09,0.49}{##1}}}
\@namedef{PY@tok@sa}{\def\PY@tc##1{\textcolor[rgb]{0.73,0.13,0.13}{##1}}}
\@namedef{PY@tok@sb}{\def\PY@tc##1{\textcolor[rgb]{0.73,0.13,0.13}{##1}}}
\@namedef{PY@tok@sc}{\def\PY@tc##1{\textcolor[rgb]{0.73,0.13,0.13}{##1}}}
\@namedef{PY@tok@dl}{\def\PY@tc##1{\textcolor[rgb]{0.73,0.13,0.13}{##1}}}
\@namedef{PY@tok@s2}{\def\PY@tc##1{\textcolor[rgb]{0.73,0.13,0.13}{##1}}}
\@namedef{PY@tok@sh}{\def\PY@tc##1{\textcolor[rgb]{0.73,0.13,0.13}{##1}}}
\@namedef{PY@tok@s1}{\def\PY@tc##1{\textcolor[rgb]{0.73,0.13,0.13}{##1}}}
\@namedef{PY@tok@mb}{\def\PY@tc##1{\textcolor[rgb]{0.40,0.40,0.40}{##1}}}
\@namedef{PY@tok@mf}{\def\PY@tc##1{\textcolor[rgb]{0.40,0.40,0.40}{##1}}}
\@namedef{PY@tok@mh}{\def\PY@tc##1{\textcolor[rgb]{0.40,0.40,0.40}{##1}}}
\@namedef{PY@tok@mi}{\def\PY@tc##1{\textcolor[rgb]{0.40,0.40,0.40}{##1}}}
\@namedef{PY@tok@il}{\def\PY@tc##1{\textcolor[rgb]{0.40,0.40,0.40}{##1}}}
\@namedef{PY@tok@mo}{\def\PY@tc##1{\textcolor[rgb]{0.40,0.40,0.40}{##1}}}
\@namedef{PY@tok@ch}{\let\PY@it=\textit\def\PY@tc##1{\textcolor[rgb]{0.24,0.48,0.48}{##1}}}
\@namedef{PY@tok@cm}{\let\PY@it=\textit\def\PY@tc##1{\textcolor[rgb]{0.24,0.48,0.48}{##1}}}
\@namedef{PY@tok@cpf}{\let\PY@it=\textit\def\PY@tc##1{\textcolor[rgb]{0.24,0.48,0.48}{##1}}}
\@namedef{PY@tok@c1}{\let\PY@it=\textit\def\PY@tc##1{\textcolor[rgb]{0.24,0.48,0.48}{##1}}}
\@namedef{PY@tok@cs}{\let\PY@it=\textit\def\PY@tc##1{\textcolor[rgb]{0.24,0.48,0.48}{##1}}}

\def\PYZbs{\char`\\}
\def\PYZus{\char`\_}
\def\PYZob{\char`\{}
\def\PYZcb{\char`\}}
\def\PYZca{\char`\^}
\def\PYZam{\char`\&}
\def\PYZlt{\char`\<}
\def\PYZgt{\char`\>}
\def\PYZsh{\char`\#}
\def\PYZpc{\char`\%}
\def\PYZdl{\char`\$}
\def\PYZhy{\char`\-}
\def\PYZsq{\char`\'}
\def\PYZdq{\char`\"}
\def\PYZti{\char`\~}
% for compatibility with earlier versions
\def\PYZat{@}
\def\PYZlb{[}
\def\PYZrb{]}
\makeatother


    % For linebreaks inside Verbatim environment from package fancyvrb. 
    \makeatletter
        \newbox\Wrappedcontinuationbox 
        \newbox\Wrappedvisiblespacebox 
        \newcommand*\Wrappedvisiblespace {\textcolor{red}{\textvisiblespace}} 
        \newcommand*\Wrappedcontinuationsymbol {\textcolor{red}{\llap{\tiny$\m@th\hookrightarrow$}}} 
        \newcommand*\Wrappedcontinuationindent {3ex } 
        \newcommand*\Wrappedafterbreak {\kern\Wrappedcontinuationindent\copy\Wrappedcontinuationbox} 
        % Take advantage of the already applied Pygments mark-up to insert 
        % potential linebreaks for TeX processing. 
        %        {, <, #, %, $, ' and ": go to next line. 
        %        _, }, ^, &, >, - and ~: stay at end of broken line. 
        % Use of \textquotesingle for straight quote. 
        \newcommand*\Wrappedbreaksatspecials {% 
            \def\PYGZus{\discretionary{\char`\_}{\Wrappedafterbreak}{\char`\_}}% 
            \def\PYGZob{\discretionary{}{\Wrappedafterbreak\char`\{}{\char`\{}}% 
            \def\PYGZcb{\discretionary{\char`\}}{\Wrappedafterbreak}{\char`\}}}% 
            \def\PYGZca{\discretionary{\char`\^}{\Wrappedafterbreak}{\char`\^}}% 
            \def\PYGZam{\discretionary{\char`\&}{\Wrappedafterbreak}{\char`\&}}% 
            \def\PYGZlt{\discretionary{}{\Wrappedafterbreak\char`\<}{\char`\<}}% 
            \def\PYGZgt{\discretionary{\char`\>}{\Wrappedafterbreak}{\char`\>}}% 
            \def\PYGZsh{\discretionary{}{\Wrappedafterbreak\char`\#}{\char`\#}}% 
            \def\PYGZpc{\discretionary{}{\Wrappedafterbreak\char`\%}{\char`\%}}% 
            \def\PYGZdl{\discretionary{}{\Wrappedafterbreak\char`\$}{\char`\$}}% 
            \def\PYGZhy{\discretionary{\char`\-}{\Wrappedafterbreak}{\char`\-}}% 
            \def\PYGZsq{\discretionary{}{\Wrappedafterbreak\textquotesingle}{\textquotesingle}}% 
            \def\PYGZdq{\discretionary{}{\Wrappedafterbreak\char`\"}{\char`\"}}% 
            \def\PYGZti{\discretionary{\char`\~}{\Wrappedafterbreak}{\char`\~}}% 
        } 
        % Some characters . , ; ? ! / are not pygmentized. 
        % This macro makes them "active" and they will insert potential linebreaks 
        \newcommand*\Wrappedbreaksatpunct {% 
            \lccode`\~`\.\lowercase{\def~}{\discretionary{\hbox{\char`\.}}{\Wrappedafterbreak}{\hbox{\char`\.}}}% 
            \lccode`\~`\,\lowercase{\def~}{\discretionary{\hbox{\char`\,}}{\Wrappedafterbreak}{\hbox{\char`\,}}}% 
            \lccode`\~`\;\lowercase{\def~}{\discretionary{\hbox{\char`\;}}{\Wrappedafterbreak}{\hbox{\char`\;}}}% 
            \lccode`\~`\:\lowercase{\def~}{\discretionary{\hbox{\char`\:}}{\Wrappedafterbreak}{\hbox{\char`\:}}}% 
            \lccode`\~`\?\lowercase{\def~}{\discretionary{\hbox{\char`\?}}{\Wrappedafterbreak}{\hbox{\char`\?}}}% 
            \lccode`\~`\!\lowercase{\def~}{\discretionary{\hbox{\char`\!}}{\Wrappedafterbreak}{\hbox{\char`\!}}}% 
            \lccode`\~`\/\lowercase{\def~}{\discretionary{\hbox{\char`\/}}{\Wrappedafterbreak}{\hbox{\char`\/}}}% 
            \catcode`\.\active
            \catcode`\,\active 
            \catcode`\;\active
            \catcode`\:\active
            \catcode`\?\active
            \catcode`\!\active
            \catcode`\/\active 
            \lccode`\~`\~ 	
        }
    \makeatother

    \let\OriginalVerbatim=\Verbatim
    \makeatletter
    \renewcommand{\Verbatim}[1][1]{%
        %\parskip\z@skip
        \sbox\Wrappedcontinuationbox {\Wrappedcontinuationsymbol}%
        \sbox\Wrappedvisiblespacebox {\FV@SetupFont\Wrappedvisiblespace}%
        \def\FancyVerbFormatLine ##1{\hsize\linewidth
            \vtop{\raggedright\hyphenpenalty\z@\exhyphenpenalty\z@
                \doublehyphendemerits\z@\finalhyphendemerits\z@
                \strut ##1\strut}%
        }%
        % If the linebreak is at a space, the latter will be displayed as visible
        % space at end of first line, and a continuation symbol starts next line.
        % Stretch/shrink are however usually zero for typewriter font.
        \def\FV@Space {%
            \nobreak\hskip\z@ plus\fontdimen3\font minus\fontdimen4\font
            \discretionary{\copy\Wrappedvisiblespacebox}{\Wrappedafterbreak}
            {\kern\fontdimen2\font}%
        }%
        
        % Allow breaks at special characters using \PYG... macros.
        \Wrappedbreaksatspecials
        % Breaks at punctuation characters . , ; ? ! and / need catcode=\active 	
        \OriginalVerbatim[#1,codes*=\Wrappedbreaksatpunct]%
    }
    \makeatother

    % Exact colors from NB
    \definecolor{incolor}{HTML}{303F9F}
    \definecolor{outcolor}{HTML}{D84315}
    \definecolor{cellborder}{HTML}{CFCFCF}
    \definecolor{cellbackground}{HTML}{F7F7F7}
    
    % prompt
    \makeatletter
    \newcommand{\boxspacing}{\kern\kvtcb@left@rule\kern\kvtcb@boxsep}
    \makeatother
    \newcommand{\prompt}[4]{
        {\ttfamily\llap{{\color{#2}[#3]:\hspace{3pt}#4}}\vspace{-\baselineskip}}
    }
    

    
    % Prevent overflowing lines due to hard-to-break entities
    \sloppy 
    % Setup hyperref package
    \hypersetup{
      breaklinks=true,  % so long urls are correctly broken across lines
      colorlinks=true,
      urlcolor=urlcolor,
      linkcolor=linkcolor,
      citecolor=citecolor,
      }
    % Slightly bigger margins than the latex defaults
    
    \geometry{verbose,tmargin=1in,bmargin=1in,lmargin=1in,rmargin=1in}
    
    

\begin{document}
    
    \maketitle
    
    

    
    \includegraphics[keepaspectratio]{Redfin.jpg}

    \section{Predicting Housing Market Trends Using Data
Science}\label{predicting-housing-market-trends-using-data-science}

\textbf{Potential Insights:} This dataset provides insights into housing
market trends across different regions. I chose this dataset because
housing market data reflects economic patterns, consumer behavior, and
financial stability. By analyzing price fluctuations, sales trends, and
inventory changes, I aim to identify key indicators that predict market
shifts. The structured approach from Kaggle competitions will be applied
to extract meaningful patterns, focusing on how market dynamics evolve
over time.

Additionally, the dataset's inclusion of multiple region types allows
for comparisons between national trends and city-specific behaviors,
providing an opportunity to explore how different real estate markets
respond to economic factors. By leveraging machine learning techniques
learned from the competitions, I will analyze potential leading
indicators of price changes and market cycles.

\subsubsection{Project Scope}\label{project-scope}

\textbf{Objectives:} Identify key trends and patterns in the housing
market, predict market fluctuations, and provide actionable insights
into real estate trend

\textbf{Deliverables:} Charts, graphs, statistical reports, predictive
models, and a final research paper detailing insights and methodologies.

Milestones:\\
- Data collection and cleaning (Week 1)\\
- Exploratory Data Analysis (Weeks 2-3)\\
- Model selection and training (Weeks 4-6)\\
- Model validation and optimization (Weeks 7-8)\\
- Final report and presentation (Weeks 9-10)

Tasks:\\
- Data acquisition and preprocessing\\
- Exploratory data analysis (EDA)\\
- Feature engineering and selection\\
- Model training and validation\\
- Performance evaluation and optimization\\
- Presentation of findings

\textbf{Resources:} Python, Pandas, NumPy, Scikit-Learn, Matplotlib,
Seaborn, Jupyter Notebook

\subsubsection{Research Plan}\label{research-plan}

\textbf{Techniques and Methods:} Time-series analysis, regression
models, machine learning techniques, and cross-validation methods.
Application: Apply time-series forecasting models to predict housing
price trends, use feature engineering techniques to identify key drivers
of price fluctuations, and validate results using robust evaluation
metrics.

\textbf{Hypothesis} - \textbf{Hypothesis Statement:} Housing market
trends can be predicted by analyzing historical price fluctuations,
inventory changes, and sales patterns. By applying machine learning
techniques, we can identify early indicators of market shifts, allowing
for better-informed real estate investment decisions.

    \section{Imports}\label{imports}

    \begin{tcolorbox}[breakable, size=fbox, boxrule=1pt, pad at break*=1mm,colback=cellbackground, colframe=cellborder]
\prompt{In}{incolor}{1}{\boxspacing}
\begin{Verbatim}[commandchars=\\\{\}]
\PY{c+c1}{\PYZsh{} my imports}
\PY{k+kn}{import} \PY{n+nn}{wrangle} \PY{k}{as} \PY{n+nn}{w}
\PY{k+kn}{import} \PY{n+nn}{explore} \PY{k}{as} \PY{n+nn}{e}
\PY{c+c1}{\PYZsh{}import modeling as m}

\PY{c+c1}{\PYZsh{} ignore warnings}
\PY{k+kn}{import} \PY{n+nn}{warnings}
\PY{n}{warnings}\PY{o}{.}\PY{n}{filterwarnings}\PY{p}{(}\PY{l+s+s2}{\PYZdq{}}\PY{l+s+s2}{ignore}\PY{l+s+s2}{\PYZdq{}}\PY{p}{)}
\end{Verbatim}
\end{tcolorbox}

    \begin{Verbatim}[commandchars=\\\{\}]
imports loaded successfully, awaiting commands{\ldots}
    \end{Verbatim}

    \subsection{\# Wrangle}\label{wrangle}

\begin{itemize}
\tightlist
\item
  Raw Data acquired from Redfin
\item
  1099 rows x 20 columns before cleaning
\item
  Data acquired from Google BigQuery after cleaning
\item
  1099 rows x 20 columns after cleaning
\end{itemize}

    \subsection{\#\# Prepare}\label{prepare}

\begin{itemize}
\tightlist
\item
  Clean the data

  \begin{itemize}
  \tightlist
  \item
    Rename columns
  \item
    Remove nulls
  \item
    Remove extra characters
  \item
    Remove extra spaces
  \item
    Fixed data type
  \end{itemize}
\end{itemize}

    \begin{tcolorbox}[breakable, size=fbox, boxrule=1pt, pad at break*=1mm,colback=cellbackground, colframe=cellborder]
\prompt{In}{incolor}{2}{\boxspacing}
\begin{Verbatim}[commandchars=\\\{\}]
\PY{n}{redfin\PYZus{}df} \PY{o}{=} \PY{n}{w}\PY{o}{.}\PY{n}{check\PYZus{}file\PYZus{}exists\PYZus{}gbq}\PY{p}{(}\PY{l+s+s2}{\PYZdq{}}\PY{l+s+s2}{data.csv}\PY{l+s+s2}{\PYZdq{}}\PY{p}{,} \PY{l+s+s2}{\PYZdq{}}\PY{l+s+s2}{service\PYZhy{}account\PYZhy{}key.json}\PY{l+s+s2}{\PYZdq{}}\PY{p}{,} 
                                    \PY{l+s+s2}{\PYZdq{}}\PY{l+s+s2}{iu\PYZhy{}thesis\PYZhy{}project.Redfin\PYZus{}Monthly\PYZus{}Housing\PYZus{}Market\PYZus{}Data.Redfin}\PY{l+s+s2}{\PYZdq{}}\PY{p}{)}
\end{Verbatim}
\end{tcolorbox}

    \begin{Verbatim}[commandchars=\\\{\}]
✅ CSV file 'data.csv' found locally. Checking data quality{\ldots}
Detected encoding: ascii
✅ Data appears to be clean.
✅ Data is already clean. No need to clean or upload.
    \end{Verbatim}

    \subsection{\#\#\# A Brief Look At The
Data}\label{a-brief-look-at-the-data}

    \begin{tcolorbox}[breakable, size=fbox, boxrule=1pt, pad at break*=1mm,colback=cellbackground, colframe=cellborder]
\prompt{In}{incolor}{3}{\boxspacing}
\begin{Verbatim}[commandchars=\\\{\}]
\PY{n}{redfin\PYZus{}df}\PY{o}{.}\PY{n}{head}\PY{p}{(}\PY{p}{)}
\end{Verbatim}
\end{tcolorbox}

            \begin{tcolorbox}[breakable, size=fbox, boxrule=.5pt, pad at break*=1mm, opacityfill=0]
\prompt{Out}{outcolor}{3}{\boxspacing}
\begin{Verbatim}[commandchars=\\\{\}]
      region       month\_of\_period\_end  median\_sale\_price  \textbackslash{}
0   National 2022-07-01 00:00:00+00:00             413000
1   National 2022-08-01 00:00:00+00:00             407000
2   National 2022-06-01 00:00:00+00:00             429000
3   National 2021-08-01 00:00:00+00:00             381000
4   National 2021-09-01 00:00:00+00:00             377000

   median\_sale\_price\_mom  median\_sale\_price\_yoy  homes\_sold  homes\_sold\_mom  \textbackslash{}
0                   -3.6                    7.4      530476           -14.4
1                   -1.5                    6.9      559150             5.4
2                   -0.7                   10.8      619857             4.1
3                   -1.1                   16.2      672612            -1.0
4                   -1.0                   13.9      638337            -5.1

   homes\_sold\_yoy  new\_listings  new\_listings\_mom  new\_listings\_yoy  \textbackslash{}
0           -21.9        705215             -14.3             -11.9
1           -16.9        622537             -11.7             -13.7
2           -15.0        822855               6.4               1.0
3             0.1        721768              -9.8               0.0
4            -4.7        675008              -6.5              -4.5

   inventory  inventory\_mom  inventory\_yoy  days\_on\_market  \textbackslash{}
0    1240506           10.2           27.1              21
1    1225537           -1.2           26.7              27
2    1125321           20.7           28.1              18
3     967307           -0.9          -18.2              17
4     969546            0.2          -16.1              19

   days\_on\_market\_mom  days\_on\_market\_yoy  average\_sale\_to\_list  \textbackslash{}
0                   3                   6                 101.0
1                   5                  10                  99.9
2                   1                   3                 102.2
3                   2                 -15                 101.6
4                   2                 -11                 101.0

   average\_sale\_to\_list\_mom  average\_sale\_to\_list\_yoy
0                      -1.2                      -1.2
1                      -1.1                      -1.7
2                      -0.8                      -0.3
3                      -0.6                       2.4
4                      -0.6                       1.7
\end{Verbatim}
\end{tcolorbox}
        
    \subsection{\#\#\# A Summary Of The Data}\label{a-summary-of-the-data}

    \begin{tcolorbox}[breakable, size=fbox, boxrule=1pt, pad at break*=1mm,colback=cellbackground, colframe=cellborder]
\prompt{In}{incolor}{4}{\boxspacing}
\begin{Verbatim}[commandchars=\\\{\}]
\PY{n}{e}\PY{o}{.}\PY{n}{data\PYZus{}summary}\PY{p}{(}\PY{n}{redfin\PYZus{}df}\PY{p}{)}
\end{Verbatim}
\end{tcolorbox}

    \begin{Verbatim}[commandchars=\\\{\}]
data shape: (1099, 20)
    \end{Verbatim}

            \begin{tcolorbox}[breakable, size=fbox, boxrule=.5pt, pad at break*=1mm, opacityfill=0]
\prompt{Out}{outcolor}{4}{\boxspacing}
\begin{Verbatim}[commandchars=\\\{\}]
                                    data type  \#missing  \%missing  \#unique  \textbackslash{}
region                                 object         0       0.0        7
month\_of\_period\_end       datetime64[ns, UTC]         0       0.0      157
median\_sale\_price                       int64         0       0.0      416
median\_sale\_price\_mom                 float64         0       0.0      186
median\_sale\_price\_yoy                 float64         0       0.0      251
homes\_sold                              int64         0       0.0     1036
homes\_sold\_mom                        float64         0       0.0      556
homes\_sold\_yoy                        float64         0       0.0      490
new\_listings                            int64         0       0.0     1062
new\_listings\_mom                      float64         0       0.0      634
new\_listings\_yoy                      float64         0       0.0      462
inventory                               int64         0       0.0     1084
inventory\_mom                         float64         0       0.0      367
inventory\_yoy                         float64         0       0.0      558
days\_on\_market                          int64         0       0.0      100
days\_on\_market\_mom                      int64         0       0.0       45
days\_on\_market\_yoy                      int64         0       0.0       86
average\_sale\_to\_list                  float64         0       0.0      130
average\_sale\_to\_list\_mom              float64         0       0.0       43
average\_sale\_to\_list\_yoy              float64         0       0.0       99

                           count                                 mean  \textbackslash{}
region                      1099                                  NaN
month\_of\_period\_end         1099  2018-07-01 19:52:21.401273856+00:00
median\_sale\_price         1099.0                        407515.013649
median\_sale\_price\_mom     1099.0                             0.610464
median\_sale\_price\_yoy     1099.0                             6.849045
homes\_sold                1099.0                         73960.789809
homes\_sold\_mom            1099.0                              1.37152
homes\_sold\_yoy            1099.0                              1.94222
new\_listings              1099.0                         89187.496815
new\_listings\_mom          1099.0                             3.776797
new\_listings\_yoy          1099.0                              0.10919
inventory                 1099.0                        224060.347589
inventory\_mom             1099.0                            -0.101001
inventory\_yoy             1099.0                            -4.468062
days\_on\_market            1099.0                            38.658781
days\_on\_market\_mom        1099.0                            -0.213831
days\_on\_market\_yoy        1099.0                            -3.524113
average\_sale\_to\_list      1099.0                            99.071338
average\_sale\_to\_list\_mom  1099.0                             0.020291
average\_sale\_to\_list\_yoy  1099.0                             0.301911

                                    std                        min  \textbackslash{}
region                              NaN                        NaN
month\_of\_period\_end                 NaN  2012-01-01 00:00:00+00:00
median\_sale\_price          189404.64317                   134000.0
median\_sale\_price\_mom          3.657105                      -12.8
median\_sale\_price\_yoy           5.72749                      -10.7
homes\_sold                174049.251215                      974.0
homes\_sold\_mom                18.047561                      -40.6
homes\_sold\_yoy                17.349035                      -56.9
new\_listings              210707.155307                     1010.0
new\_listings\_mom              29.206833                      -59.5
new\_listings\_yoy              17.261469                      -70.1
inventory                 541161.876914                     1057.0
inventory\_mom                  9.783868                      -39.3
inventory\_yoy                 23.895725                      -60.5
days\_on\_market                21.157081                        5.0
days\_on\_market\_mom             6.726423                      -49.0
days\_on\_market\_yoy             12.70981                      -69.0
average\_sale\_to\_list           2.213431                       92.0
average\_sale\_to\_list\_mom       0.510114                       -4.1
average\_sale\_to\_list\_yoy       1.418228                      -10.6

                                                25\%  \textbackslash{}
region                                          NaN
month\_of\_period\_end       2015-04-01 00:00:00+00:00
median\_sale\_price                          256500.0
median\_sale\_price\_mom                          -1.8
median\_sale\_price\_yoy                          3.35
homes\_sold                                   3283.0
homes\_sold\_mom                                -10.7
homes\_sold\_yoy                                 -6.5
new\_listings                                 3613.0
new\_listings\_mom                              -13.0
new\_listings\_yoy                               -6.8
inventory                                    7138.0
inventory\_mom                                 -4.45
inventory\_yoy                                 -17.4
days\_on\_market                                 23.0
days\_on\_market\_mom                             -2.0
days\_on\_market\_yoy                             -8.0
average\_sale\_to\_list                           97.7
average\_sale\_to\_list\_mom                       -0.2
average\_sale\_to\_list\_yoy                       -0.1

                                                50\%  \textbackslash{}
region                                          NaN
month\_of\_period\_end       2018-07-01 00:00:00+00:00
median\_sale\_price                          371000.0
median\_sale\_price\_mom                           0.2
median\_sale\_price\_yoy                           6.5
homes\_sold                                   5108.0
homes\_sold\_mom                                  0.2
homes\_sold\_yoy                                  1.8
new\_listings                                 6070.0
new\_listings\_mom                               -2.5
new\_listings\_yoy                                0.2
inventory                                   12397.0
inventory\_mom                                   0.4
inventory\_yoy                                  -7.3
days\_on\_market                                 36.0
days\_on\_market\_mom                              1.0
days\_on\_market\_yoy                             -2.0
average\_sale\_to\_list                           99.1
average\_sale\_to\_list\_mom                       -0.1
average\_sale\_to\_list\_yoy                        0.3

                                                75\%                        max
region                                          NaN                        NaN
month\_of\_period\_end       2021-10-01 00:00:00+00:00  2025-01-01 00:00:00+00:00
median\_sale\_price                          525000.0                   925000.0
median\_sale\_price\_mom                           2.8                       16.4
median\_sale\_price\_yoy                           9.7                       31.8
homes\_sold                                   7373.0                   729414.0
homes\_sold\_mom                                11.55                       63.4
homes\_sold\_yoy                                 10.7                      127.6
new\_listings                                 9152.0                   844234.0
new\_listings\_mom                              14.05                      150.3
new\_listings\_yoy                               6.75                      257.9
inventory                                   22078.0                  2181834.0
inventory\_mom                                  4.55                       49.4
inventory\_yoy                                   4.9                      230.0
days\_on\_market                                 49.0                      145.0
days\_on\_market\_mom                              4.0                       19.0
days\_on\_market\_yoy                              3.0                       43.0
average\_sale\_to\_list                          100.2                      110.7
average\_sale\_to\_list\_mom                        0.3                        3.6
average\_sale\_to\_list\_yoy                        0.9                        7.2
\end{Verbatim}
\end{tcolorbox}
        
    \begin{tcolorbox}[breakable, size=fbox, boxrule=1pt, pad at break*=1mm,colback=cellbackground, colframe=cellborder]
\prompt{In}{incolor}{5}{\boxspacing}
\begin{Verbatim}[commandchars=\\\{\}]
\PY{n}{redfin\PYZus{}df}\PY{o}{.}\PY{n}{info}\PY{p}{(}\PY{p}{)}
\end{Verbatim}
\end{tcolorbox}

    \begin{Verbatim}[commandchars=\\\{\}]
<class 'pandas.core.frame.DataFrame'>
RangeIndex: 1099 entries, 0 to 1098
Data columns (total 20 columns):
 \#   Column                    Non-Null Count  Dtype
---  ------                    --------------  -----
 0   region                    1099 non-null   object
 1   month\_of\_period\_end       1099 non-null   datetime64[ns, UTC]
 2   median\_sale\_price         1099 non-null   int64
 3   median\_sale\_price\_mom     1099 non-null   float64
 4   median\_sale\_price\_yoy     1099 non-null   float64
 5   homes\_sold                1099 non-null   int64
 6   homes\_sold\_mom            1099 non-null   float64
 7   homes\_sold\_yoy            1099 non-null   float64
 8   new\_listings              1099 non-null   int64
 9   new\_listings\_mom          1099 non-null   float64
 10  new\_listings\_yoy          1099 non-null   float64
 11  inventory                 1099 non-null   int64
 12  inventory\_mom             1099 non-null   float64
 13  inventory\_yoy             1099 non-null   float64
 14  days\_on\_market            1099 non-null   int64
 15  days\_on\_market\_mom        1099 non-null   int64
 16  days\_on\_market\_yoy        1099 non-null   int64
 17  average\_sale\_to\_list      1099 non-null   float64
 18  average\_sale\_to\_list\_mom  1099 non-null   float64
 19  average\_sale\_to\_list\_yoy  1099 non-null   float64
dtypes: datetime64[ns, UTC](1), float64(11), int64(7), object(1)
memory usage: 171.8+ KB
    \end{Verbatim}

    \begin{center}\rule{0.5\linewidth}{0.5pt}\end{center}

\section{Exploration}\label{exploration}

    \begin{tcolorbox}[breakable, size=fbox, boxrule=1pt, pad at break*=1mm,colback=cellbackground, colframe=cellborder]
\prompt{In}{incolor}{6}{\boxspacing}
\begin{Verbatim}[commandchars=\\\{\}]
\PY{c+c1}{\PYZsh{} visualized your data}
\PY{k+kn}{import} \PY{n+nn}{matplotlib}\PY{n+nn}{.}\PY{n+nn}{pyplot} \PY{k}{as} \PY{n+nn}{plt}
\PY{k+kn}{import} \PY{n+nn}{seaborn} \PY{k}{as} \PY{n+nn}{sns}

\PY{c+c1}{\PYZsh{} to view a bit closer the pairplot double click on the plot}
\PY{n}{sns}\PY{o}{.}\PY{n}{pairplot}\PY{p}{(}\PY{n}{redfin\PYZus{}df}\PY{p}{)}
\end{Verbatim}
\end{tcolorbox}

            \begin{tcolorbox}[breakable, size=fbox, boxrule=.5pt, pad at break*=1mm, opacityfill=0]
\prompt{Out}{outcolor}{6}{\boxspacing}
\begin{Verbatim}[commandchars=\\\{\}]
<seaborn.axisgrid.PairGrid at 0x332aa2130>
\end{Verbatim}
\end{tcolorbox}
        
    \begin{center}
    \adjustimage{max size={0.9\linewidth}{0.9\paperheight}}{thesis_project_files/thesis_project_13_1.png}
    \end{center}
    { \hspace*{\fill} \\}
    
    \begin{center}\rule{0.5\linewidth}{0.5pt}\end{center}

\paragraph{What is the distribution of the
Target?}\label{what-is-the-distribution-of-the-target}

\begin{itemize}
\tightlist
\item
  Note: this question is more of a look at the data, not a stats test
  question.
\end{itemize}

    \begin{tcolorbox}[breakable, size=fbox, boxrule=1pt, pad at break*=1mm,colback=cellbackground, colframe=cellborder]
\prompt{In}{incolor}{7}{\boxspacing}
\begin{Verbatim}[commandchars=\\\{\}]
\PY{c+c1}{\PYZsh{} get a histplot of quality}
\PY{n}{sns}\PY{o}{.}\PY{n}{histplot}\PY{p}{(}\PY{n}{redfin\PYZus{}df}\PY{o}{.}\PY{n}{median\PYZus{}sale\PYZus{}price}\PY{p}{)}
\PY{n}{plt}\PY{o}{.}\PY{n}{show}\PY{p}{(}\PY{p}{)}
\end{Verbatim}
\end{tcolorbox}

    \begin{center}
    \adjustimage{max size={0.9\linewidth}{0.9\paperheight}}{thesis_project_files/thesis_project_15_0.png}
    \end{center}
    { \hspace*{\fill} \\}
    
    \paragraph{\texorpdfstring{Target Justification:
\textbf{median\_sale\_price}}{Target Justification: median\_sale\_price}}\label{target-justification-median_sale_price}

\begin{enumerate}
\def\labelenumi{\arabic{enumi}.}
\tightlist
\item
  Core Predictive Goal

  \begin{itemize}
  \tightlist
  \item
    Both the Project Proposal (focused on housing market trend
    prediction) and Literature Review (highlighting house-price
    forecasting as critical) point to home prices as the principal
    outcome of interest.
  \item
    As emphasized in Assignment 2, ``Accurately predicting real estate
    house prices is critical for financial institutions, investors, and
    policymakers.''
  \end{itemize}
\item
  Data Alignment

  \begin{itemize}
  \tightlist
  \item
    The Redfin dataset explicitly provides median\_sale\_price (along
    with month-over-month and year-over-year variations).
  \item
    In real estate forecasting literature, the median sale price is a
    standard measure, serving as a robust, central indicator of market
    conditions.
  \end{itemize}
\item
  Common Industry Metric

  \begin{itemize}
  \tightlist
  \item
    Forecasting median sale price is standard practice in both industry
    and academic settings. It allows analysis of market health and trend
    over time, region by region.
  \end{itemize}
\end{enumerate}

    Histogram Insights

Looking at the histogram of median\_sale\_price: - The distribution
peaks roughly in the \textbackslash\$200,000 - \textbackslash\$300,000
range, indicating that many observations cluster around these mid-range
house prices. - The tail extends toward higher price points (up to
\$900,000), reflecting that while most houses fall in the lower-to-mid
range, there are notable but fewer properties in higher price brackets.
- The shape suggests a right-skew (common in housing data), where a
larger number of transactions occur at relatively moderate price levels,
and fewer yet higher-priced properties cause the tail to stretch to the
right.

\begin{center}\rule{0.5\linewidth}{0.5pt}\end{center}

    \begin{tcolorbox}[breakable, size=fbox, boxrule=1pt, pad at break*=1mm,colback=cellbackground, colframe=cellborder]
\prompt{In}{incolor}{8}{\boxspacing}
\begin{Verbatim}[commandchars=\\\{\}]
\PY{c+c1}{\PYZsh{} Correlation heatmap for all columns}
\PY{n}{e}\PY{o}{.}\PY{n}{plot\PYZus{}correlation\PYZus{}heatmap}\PY{p}{(}\PY{n}{redfin\PYZus{}df}\PY{p}{)}

\PY{c+c1}{\PYZsh{} Correlation heatmap for a subset of columns}
\PY{n}{cols\PYZus{}of\PYZus{}interest} \PY{o}{=} \PY{p}{[}\PY{l+s+s1}{\PYZsq{}}\PY{l+s+s1}{median\PYZus{}sale\PYZus{}price}\PY{l+s+s1}{\PYZsq{}}\PY{p}{,} \PY{l+s+s1}{\PYZsq{}}\PY{l+s+s1}{inventory}\PY{l+s+s1}{\PYZsq{}}\PY{p}{,} \PY{l+s+s1}{\PYZsq{}}\PY{l+s+s1}{days\PYZus{}on\PYZus{}market}\PY{l+s+s1}{\PYZsq{}}\PY{p}{]}
\PY{n}{e}\PY{o}{.}\PY{n}{plot\PYZus{}correlation\PYZus{}heatmap}\PY{p}{(}\PY{n}{redfin\PYZus{}df}\PY{p}{,} \PY{n}{columns}\PY{o}{=}\PY{n}{cols\PYZus{}of\PYZus{}interest}\PY{p}{,} \PY{n}{figsize}\PY{o}{=}\PY{p}{(}\PY{l+m+mi}{10}\PY{p}{,}\PY{l+m+mi}{8}\PY{p}{)}\PY{p}{,} \PY{n}{cmap}\PY{o}{=}\PY{l+s+s1}{\PYZsq{}}\PY{l+s+s1}{viridis}\PY{l+s+s1}{\PYZsq{}}\PY{p}{)}
\end{Verbatim}
\end{tcolorbox}

    \begin{center}
    \adjustimage{max size={0.9\linewidth}{0.9\paperheight}}{thesis_project_files/thesis_project_18_0.png}
    \end{center}
    { \hspace*{\fill} \\}
    
    \begin{center}
    \adjustimage{max size={0.9\linewidth}{0.9\paperheight}}{thesis_project_files/thesis_project_18_1.png}
    \end{center}
    { \hspace*{\fill} \\}
    
    Correlation Matrix (heatmap) Insights

\textbf{Notable Correlation Patterns} 1. \textbf{Negative} Correlation:
median\_sale\_price vs.~days\_on\_market - Around -0.47 for the
current-month relationship - Suggests that when prices are higher,
properties tend to sell faster (fewer days on market). Conversely, in
slower markets with longer days on market, median prices are typically
lower. 2. \textbf{Positive} Correlation Between ``\_yoy'' Columns - You
may see strong positive correlations among ``year-over-year''
metrics---e.g., new\_listings\_yoy vs.~inventory\_yoy. - Indicates that
increases in new listings over last year often coincide with increases
in inventory relative to last year. 3. Moderate/Strong Relationship:
median\_sale\_price and average\_sale\_to\_list - Look for a reddish
block in the row/column intersection of median\_sale\_price and
average\_sale\_to\_list (and possibly their \_yoy versions). - A higher
sale-to-list ratio can coincide with a stronger market, often reflected
in higher median prices. 4. Month-over-Month (MoM) Correlations - Some
MoM columns (like homes\_sold\_mom or median\_sale\_price\_mom) may show
moderate correlations with each other. - This could imply, for example,
that monthly surges in sales might accompany monthly rises in price---or
vice versa. 5. Days on Market and Homes Sold - Often, a negative
correlation emerges between days\_on\_market and homes\_sold or
homes\_sold\_mom: markets with more successful sales typically see homes
spending fewer days on the market.

\textbf{Interpreting High Correlation Values} - Close to +1: Two
features move in the same direction nearly every time (e.g.,
new\_listings vs.~inventory might be highly correlated if more listings
generally lead to more inventory).\\
- Close to -1: A strong inverse relationship (e.g., higher
median\_sale\_price typically corresponds to fewer days\_on\_market).\\
- Close to 0: Little direct linear relationship. There may be other
complex or nonlinear relationships, or it may just indicate no
meaningful interaction.

\textbf{Caveats \& Next Steps} - Correlation ≠ Causation: Even if
median\_sale\_price is strongly correlated with inventory, it doesn't
confirm a direct cause-effect. External factors like interest rates or
broader economic conditions could drive both.\\
- Outliers \& Nonlinearities: Some correlations may appear weaker or
stronger if there are a few extreme outliers or nonlinear relationships.
Pair plots, scatterplots, or transformations (e.g., log scales) can help
clarify.\\
- Modeling Implications: Strongly correlated variables may lead to
multicollinearity in linear models, so we will consider advanced feature
selection or regularization (e.g., Ridge/Lasso) if that's an issue.

\begin{center}\rule{0.5\linewidth}{0.5pt}\end{center}

    \begin{tcolorbox}[breakable, size=fbox, boxrule=1pt, pad at break*=1mm,colback=cellbackground, colframe=cellborder]
\prompt{In}{incolor}{9}{\boxspacing}
\begin{Verbatim}[commandchars=\\\{\}]
\PY{n}{e}\PY{o}{.}\PY{n}{plot\PYZus{}seasonal\PYZus{}decomposition}\PY{p}{(}\PY{n}{redfin\PYZus{}df}\PY{p}{,} \PY{n}{date\PYZus{}col}\PY{o}{=}\PY{l+s+s1}{\PYZsq{}}\PY{l+s+s1}{month\PYZus{}of\PYZus{}period\PYZus{}end}\PY{l+s+s1}{\PYZsq{}}\PY{p}{,} \PY{n}{target\PYZus{}col}\PY{o}{=}\PY{l+s+s1}{\PYZsq{}}\PY{l+s+s1}{median\PYZus{}sale\PYZus{}price}\PY{l+s+s1}{\PYZsq{}}\PY{p}{,} \PY{n}{model}\PY{o}{=}\PY{l+s+s1}{\PYZsq{}}\PY{l+s+s1}{multiplicative}\PY{l+s+s1}{\PYZsq{}}\PY{p}{,} \PY{n}{period}\PY{o}{=}\PY{l+m+mi}{12}\PY{p}{)}
\end{Verbatim}
\end{tcolorbox}

    \begin{center}
    \adjustimage{max size={0.9\linewidth}{0.9\paperheight}}{thesis_project_files/thesis_project_20_0.png}
    \end{center}
    { \hspace*{\fill} \\}
    
    Seasonal Decomposition Insights

\begin{enumerate}
\def\labelenumi{\arabic{enumi}.}
\tightlist
\item
  Observed

  \begin{itemize}
  \tightlist
  \item
    The top series shows a long-term increase in median sale prices from
    around 2012 to 2024.
  \item
    Notice the monthly fluctuations---some months are higher or lower
    than the smoothed trend, consistent with typical real estate
    seasonality.
  \end{itemize}
\item
  Trend

  \begin{itemize}
  \tightlist
  \item
    The middle-top subplot reveals a steady upward slope from roughly
    \textbackslash\$200k to \textbackslash\$500k+ over the last decade,
    reflecting long-term market growth.
  \item
    After 2022--2023, the trend appears to plateau or slow, suggesting
    potential cooling or stabilization in prices.
  \end{itemize}
\item
  Seasonal

  \begin{itemize}
  \tightlist
  \item
    Because the model is multiplicative, the seasonal factor hovers
    around 1.0---meaning monthly prices are typically within a few
    percentage points above or below the overall trend.
  \item
    The real estate market often shows seasonal peaks (e.g.,
    spring/summer) and troughs (winter). In a multiplicative model,
    these bumps are a percentage of the underlying trend, rather than a
    fixed amount.
  \end{itemize}
\item
  Residual

  \begin{itemize}
  \tightlist
  \item
    The bottom subplot shows what's left after removing the trend and
    seasonal components.
  \item
    If the residual looks mostly random, that indicates the
    trend+seasonality capture the main structure of the series.
    Occasional spikes or dips might point to external shocks (e.g.,
    policy changes or unexpected events).
  \end{itemize}
\end{enumerate}

\begin{center}\rule{0.5\linewidth}{0.5pt}\end{center}

    \begin{tcolorbox}[breakable, size=fbox, boxrule=1pt, pad at break*=1mm,colback=cellbackground, colframe=cellborder]
\prompt{In}{incolor}{10}{\boxspacing}
\begin{Verbatim}[commandchars=\\\{\}]
\PY{c+c1}{\PYZsh{} Identifying numerical and categorical columns}
\PY{n}{numerical\PYZus{}cols} \PY{o}{=} \PY{n}{redfin\PYZus{}df}\PY{o}{.}\PY{n}{select\PYZus{}dtypes}\PY{p}{(}\PY{n}{include}\PY{o}{=}\PY{p}{[}\PY{l+s+s1}{\PYZsq{}}\PY{l+s+s1}{int64}\PY{l+s+s1}{\PYZsq{}}\PY{p}{,} \PY{l+s+s1}{\PYZsq{}}\PY{l+s+s1}{float64}\PY{l+s+s1}{\PYZsq{}}\PY{p}{]}\PY{p}{)}\PY{o}{.}\PY{n}{columns}
\PY{n}{categorical\PYZus{}cols} \PY{o}{=} \PY{n}{redfin\PYZus{}df}\PY{o}{.}\PY{n}{select\PYZus{}dtypes}\PY{p}{(}\PY{n}{include}\PY{o}{=}\PY{p}{[}\PY{l+s+s1}{\PYZsq{}}\PY{l+s+s1}{object}\PY{l+s+s1}{\PYZsq{}}\PY{p}{,} \PY{l+s+s1}{\PYZsq{}}\PY{l+s+s1}{bool}\PY{l+s+s1}{\PYZsq{}}\PY{p}{]}\PY{p}{)}\PY{o}{.}\PY{n}{columns}
\end{Verbatim}
\end{tcolorbox}

    \begin{tcolorbox}[breakable, size=fbox, boxrule=1pt, pad at break*=1mm,colback=cellbackground, colframe=cellborder]
\prompt{In}{incolor}{11}{\boxspacing}
\begin{Verbatim}[commandchars=\\\{\}]
\PY{n}{e}\PY{o}{.}\PY{n}{bivariate\PYZus{}boxplot}\PY{p}{(}\PY{n}{redfin\PYZus{}df}\PY{p}{,} \PY{n}{numerical\PYZus{}cols}\PY{p}{,} \PY{n}{categorical\PYZus{}cols}\PY{p}{)}
\end{Verbatim}
\end{tcolorbox}

    \begin{center}
    \adjustimage{max size={0.9\linewidth}{0.9\paperheight}}{thesis_project_files/thesis_project_23_0.png}
    \end{center}
    { \hspace*{\fill} \\}
    
    \begin{center}
    \adjustimage{max size={0.9\linewidth}{0.9\paperheight}}{thesis_project_files/thesis_project_23_1.png}
    \end{center}
    { \hspace*{\fill} \\}
    
    \begin{center}
    \adjustimage{max size={0.9\linewidth}{0.9\paperheight}}{thesis_project_files/thesis_project_23_2.png}
    \end{center}
    { \hspace*{\fill} \\}
    
    \begin{center}
    \adjustimage{max size={0.9\linewidth}{0.9\paperheight}}{thesis_project_files/thesis_project_23_3.png}
    \end{center}
    { \hspace*{\fill} \\}
    
    \begin{center}
    \adjustimage{max size={0.9\linewidth}{0.9\paperheight}}{thesis_project_files/thesis_project_23_4.png}
    \end{center}
    { \hspace*{\fill} \\}
    
    \begin{center}
    \adjustimage{max size={0.9\linewidth}{0.9\paperheight}}{thesis_project_files/thesis_project_23_5.png}
    \end{center}
    { \hspace*{\fill} \\}
    
    \begin{center}
    \adjustimage{max size={0.9\linewidth}{0.9\paperheight}}{thesis_project_files/thesis_project_23_6.png}
    \end{center}
    { \hspace*{\fill} \\}
    
    \begin{center}
    \adjustimage{max size={0.9\linewidth}{0.9\paperheight}}{thesis_project_files/thesis_project_23_7.png}
    \end{center}
    { \hspace*{\fill} \\}
    
    \begin{center}
    \adjustimage{max size={0.9\linewidth}{0.9\paperheight}}{thesis_project_files/thesis_project_23_8.png}
    \end{center}
    { \hspace*{\fill} \\}
    
    \begin{center}
    \adjustimage{max size={0.9\linewidth}{0.9\paperheight}}{thesis_project_files/thesis_project_23_9.png}
    \end{center}
    { \hspace*{\fill} \\}
    
    \begin{center}
    \adjustimage{max size={0.9\linewidth}{0.9\paperheight}}{thesis_project_files/thesis_project_23_10.png}
    \end{center}
    { \hspace*{\fill} \\}
    
    \begin{center}
    \adjustimage{max size={0.9\linewidth}{0.9\paperheight}}{thesis_project_files/thesis_project_23_11.png}
    \end{center}
    { \hspace*{\fill} \\}
    
    \begin{center}
    \adjustimage{max size={0.9\linewidth}{0.9\paperheight}}{thesis_project_files/thesis_project_23_12.png}
    \end{center}
    { \hspace*{\fill} \\}
    
    \begin{center}
    \adjustimage{max size={0.9\linewidth}{0.9\paperheight}}{thesis_project_files/thesis_project_23_13.png}
    \end{center}
    { \hspace*{\fill} \\}
    
    \begin{center}
    \adjustimage{max size={0.9\linewidth}{0.9\paperheight}}{thesis_project_files/thesis_project_23_14.png}
    \end{center}
    { \hspace*{\fill} \\}
    
    \begin{center}
    \adjustimage{max size={0.9\linewidth}{0.9\paperheight}}{thesis_project_files/thesis_project_23_15.png}
    \end{center}
    { \hspace*{\fill} \\}
    
    \begin{center}
    \adjustimage{max size={0.9\linewidth}{0.9\paperheight}}{thesis_project_files/thesis_project_23_16.png}
    \end{center}
    { \hspace*{\fill} \\}
    
    \begin{center}
    \adjustimage{max size={0.9\linewidth}{0.9\paperheight}}{thesis_project_files/thesis_project_23_17.png}
    \end{center}
    { \hspace*{\fill} \\}
    
    \begin{center}
    \adjustimage{max size={0.9\linewidth}{0.9\paperheight}}{thesis_project_files/thesis_project_23_18.png}
    \end{center}
    { \hspace*{\fill} \\}
    
    BoxPlot Insights

Below is an explanation of the key boxplots I've produced and why they
matter for my project. Rather than listing every single variable, I'm
focusing on those with notable outliers or high relevance to house-price
modeling. 1. Homes Sold (Boxplot) - Notable Outliers: A few points stand
far above the upper whisker, indicating especially high sales in certain
months/regions. - Relevance: Large spikes in homes sold can influence
average prices and reflect heightened demand. Such outliers may indicate
unusually hot markets or seasonal booms (e.g., spring). 2. Days on
Market (Boxplot) - Outliers: Some properties linger on the market
significantly longer than typical, which can extend well beyond 100+
days. - Relevance: Days on market is inversely correlated with
price---when houses sell quickly, prices are often higher. The outliers
could represent unique property types, overpriced listings, or less
desirable locations. 3. Inventory (Boxplot) - Outliers: A handful of
months/regions show extremely high inventory levels, well above typical
ranges. - Relevance: Excessive inventory can apply downward pressure on
prices. Outlier months might coincide with shifts in economic conditions
or new construction spikes. 4. Median Sale Price (Boxplot) - Skew \&
High-Price Outliers: The boxplot is right-skewed; some luxury or
high-demand areas push prices well above the median. - Relevance: These
outliers are integral to your analysis, because price is the target
variable. Log transformation or robust regressions might help handle
these extremes in modeling. 5. Boxplot by price\_bin - Why Binning?:
Converting median\_sale\_price to categories (e.g., Low, Medium, High)
allows us to see how other metrics (like homes\_sold, days\_on\_market)
differ across price tiers. - Key Observations: Higher-tier price bins
often correlate with fewer days on market and slightly lower inventory,
suggesting strong demand in more expensive markets. A small subset of
extreme outliers can appear in each bin (e.g., unusually fast or slow
sales).

\textbf{Next Steps}\\
- Decide on Outlier Treatment: For modeling, should I exclude or
transform extreme values?\\
- Connect Boxplots to Time-Series: Check whether outlier months line up
with anomalies in my seasonal decomposition or other metrics (e.g.,
interest rate changes).\\
- Refine Feature Selection: Outliers can distort correlations. I will
consider robust scaling or transformations (log, sqrt) to reduce skew.

\begin{center}\rule{0.5\linewidth}{0.5pt}\end{center}

    \begin{tcolorbox}[breakable, size=fbox, boxrule=1pt, pad at break*=1mm,colback=cellbackground, colframe=cellborder]
\prompt{In}{incolor}{12}{\boxspacing}
\begin{Verbatim}[commandchars=\\\{\}]
\PY{n}{e}\PY{o}{.}\PY{n}{bivariate\PYZus{}scatterplot}\PY{p}{(}\PY{n}{redfin\PYZus{}df}\PY{p}{,} \PY{n}{numerical\PYZus{}cols}\PY{p}{,} \PY{n}{categorical\PYZus{}cols}\PY{p}{)}
\end{Verbatim}
\end{tcolorbox}

    \begin{center}
    \adjustimage{max size={0.9\linewidth}{0.9\paperheight}}{thesis_project_files/thesis_project_25_0.png}
    \end{center}
    { \hspace*{\fill} \\}
    
    \begin{center}
    \adjustimage{max size={0.9\linewidth}{0.9\paperheight}}{thesis_project_files/thesis_project_25_1.png}
    \end{center}
    { \hspace*{\fill} \\}
    
    \begin{center}
    \adjustimage{max size={0.9\linewidth}{0.9\paperheight}}{thesis_project_files/thesis_project_25_2.png}
    \end{center}
    { \hspace*{\fill} \\}
    
    \begin{center}
    \adjustimage{max size={0.9\linewidth}{0.9\paperheight}}{thesis_project_files/thesis_project_25_3.png}
    \end{center}
    { \hspace*{\fill} \\}
    
    \begin{center}
    \adjustimage{max size={0.9\linewidth}{0.9\paperheight}}{thesis_project_files/thesis_project_25_4.png}
    \end{center}
    { \hspace*{\fill} \\}
    
    \begin{center}
    \adjustimage{max size={0.9\linewidth}{0.9\paperheight}}{thesis_project_files/thesis_project_25_5.png}
    \end{center}
    { \hspace*{\fill} \\}
    
    \begin{center}
    \adjustimage{max size={0.9\linewidth}{0.9\paperheight}}{thesis_project_files/thesis_project_25_6.png}
    \end{center}
    { \hspace*{\fill} \\}
    
    \begin{center}
    \adjustimage{max size={0.9\linewidth}{0.9\paperheight}}{thesis_project_files/thesis_project_25_7.png}
    \end{center}
    { \hspace*{\fill} \\}
    
    \begin{center}
    \adjustimage{max size={0.9\linewidth}{0.9\paperheight}}{thesis_project_files/thesis_project_25_8.png}
    \end{center}
    { \hspace*{\fill} \\}
    
    \begin{center}
    \adjustimage{max size={0.9\linewidth}{0.9\paperheight}}{thesis_project_files/thesis_project_25_9.png}
    \end{center}
    { \hspace*{\fill} \\}
    
    \begin{center}
    \adjustimage{max size={0.9\linewidth}{0.9\paperheight}}{thesis_project_files/thesis_project_25_10.png}
    \end{center}
    { \hspace*{\fill} \\}
    
    \begin{center}
    \adjustimage{max size={0.9\linewidth}{0.9\paperheight}}{thesis_project_files/thesis_project_25_11.png}
    \end{center}
    { \hspace*{\fill} \\}
    
    \begin{center}
    \adjustimage{max size={0.9\linewidth}{0.9\paperheight}}{thesis_project_files/thesis_project_25_12.png}
    \end{center}
    { \hspace*{\fill} \\}
    
    \begin{center}
    \adjustimage{max size={0.9\linewidth}{0.9\paperheight}}{thesis_project_files/thesis_project_25_13.png}
    \end{center}
    { \hspace*{\fill} \\}
    
    \begin{center}
    \adjustimage{max size={0.9\linewidth}{0.9\paperheight}}{thesis_project_files/thesis_project_25_14.png}
    \end{center}
    { \hspace*{\fill} \\}
    
    \begin{center}
    \adjustimage{max size={0.9\linewidth}{0.9\paperheight}}{thesis_project_files/thesis_project_25_15.png}
    \end{center}
    { \hspace*{\fill} \\}
    
    \begin{center}
    \adjustimage{max size={0.9\linewidth}{0.9\paperheight}}{thesis_project_files/thesis_project_25_16.png}
    \end{center}
    { \hspace*{\fill} \\}
    
    ScatterPlot Insights

Below is an interpretation of the scatterplots I've generated.

\textbf{1. Overarching Themes \& Patterns}\\
1. Region-Specific Clustering - Boston (Orange) repeatedly stands out
for higher median prices
(\textasciitilde\textbackslash\$550k--\textbackslash\$600k) with short
days on market (under \textasciitilde25 days), suggesting a hot market
with sustained demand. - Chicago (Green) shows moderate prices
(\textasciitilde\textbackslash\$300k--\textbackslash\$400k) but
occasional spikes above 60 days on market, implying seasonal or economic
slowdowns. - Los Angeles (Red) often splits into two sub-markets: quick
sales at moderate-to-high price levels and pricier listings lingering
above 50 days on market, reflecting LA's diverse property segments. -
Seattle (Brown) experiences inventory surges that don't immediately drag
down prices---an early sign that, if demand weakens, a more noticeable
price dip could follow. 2. Inventory vs.~Price Lags - Across multiple
scatterplots (especially for Seattle and the National data in Blue), we
see rising inventory levels without an immediate drop in median sale
prices. - This lag effect suggests that supply growth takes time to
influence final sale prices---aligning with the idea that inventory can
be an early warning measure. 3. Sales Volumes \& Days on Market -
Washington, DC (Pink) and Philadelphia (Purple) generally hover in a
midrange of 20--40 days on market, with moderate price points. However,
DC shows a relatively high volume of sales, indicating quick condo
turnover or robust buyer interest. - When days on market does spike, the
scatterplots show it often corresponds to slower sales volume in certain
months, clarifying earlier boxplot outliers as tied to specific time
frames or local conditions. --- \textbf{2. New Insights Beyond the
Boxplots}\\
1. Identifying Hidden Anomalies - Several scatterplots reveal ``micro
outliers'' that don't appear extreme enough in boxplots to be flagged,
but stand out when comparing two variables (e.g., extremely fast sales
in a mid-price bracket for a specific region or month). - These subtle
anomalies can signal pockets of unique activity, such as sudden investor
interest or localized job growth, invisible in a univariate boxplot. 2.
Leading Indicators for Market Shifts - In LA, DC, and Seattle
especially, we see how an uptick in days on market or a surge in
inventory often foreshadows a slight flattening or gradual decrease in
prices a few months later. - This suggests a predictive dimension: if
these patterns are captured in a time-series or machine learning model,
they might serve as early warning signs for price shifts. 3. Region
vs.~National Discrepancies - The National (Blue) plots smooth out local
spikes, confirming that aggregated data can mask volatile patterns
visible only when looking at individual metros. - This underscores the
importance of regional modeling: markets like Boston or LA behave
differently than the national average, which could improve the
hypothesis about analyzing local supply-demand factors. --- \textbf{3.
Tying It All to the Hypothesis}

Hypothesis Statement: Housing market trends can be predicted by
analyzing historical price fluctuations, inventory changes, and sales
patterns. By applying machine learning techniques, we can identify early
indicators of market shifts, allowing for better-informed real estate
investment decisions.

\begin{itemize}
\tightlist
\item
  The scatterplots confirm that:

  \begin{enumerate}
  \def\labelenumi{\arabic{enumi}.}
  \tightlist
  \item
    Price, inventory, and days on market each show distinct patterns by
    region.
  \item
    Lag effects (inventory changes preceding price moves) frequently
    appear, signifying an early signal capability.
  \item
    Outliers and sub-market clusters identified here reinforce the need
    to incorporate region-specific characteristics into any predictive
    model.
  \end{enumerate}
\end{itemize}

With these visual confirmations, the data strongly supports the
hypothesis that local inventory spikes, fluctuating days on market, and
changes in sales volume can foretell upcoming price adjustments. A
machine learning approach integrating these features---especially broken
down by metro area---stands a good chance of capturing and predicting
market shifts more accurately than a one-size-fits-all national model.

\begin{center}\rule{0.5\linewidth}{0.5pt}\end{center}

    \begin{tcolorbox}[breakable, size=fbox, boxrule=1pt, pad at break*=1mm,colback=cellbackground, colframe=cellborder]
\prompt{In}{incolor}{13}{\boxspacing}
\begin{Verbatim}[commandchars=\\\{\}]
\PY{n}{redfin\PYZus{}df}\PY{o}{.}\PY{n}{groupby}\PY{p}{(}\PY{l+s+s1}{\PYZsq{}}\PY{l+s+s1}{region}\PY{l+s+s1}{\PYZsq{}}\PY{p}{)}\PY{p}{[}\PY{l+s+s1}{\PYZsq{}}\PY{l+s+s1}{price\PYZus{}bin}\PY{l+s+s1}{\PYZsq{}}\PY{p}{]}\PY{o}{.}\PY{n}{describe}\PY{p}{(}\PY{p}{)}
\end{Verbatim}
\end{tcolorbox}

            \begin{tcolorbox}[breakable, size=fbox, boxrule=.5pt, pad at break*=1mm, opacityfill=0]
\prompt{Out}{outcolor}{13}{\boxspacing}
\begin{Verbatim}[commandchars=\\\{\}]
                            count unique        top freq
region
 National                     157      3        Low   61
Boston, MA metro area         157      3       High   71
Chicago, IL metro area        157      3        Low   94
Los Angeles, CA metro area    157      3  Very High  106
Philadelphia, PA metro area   157      2        Low  119
Seattle, WA metro area        157      4  Very High   85
Washington, DC metro area     157      3       High   83
\end{Verbatim}
\end{tcolorbox}
        
    Price\_bin distribution by region Insights

Below is overview of the price\_bin distribution by region. Each region
has 157 rows of data, and the table shows which bin is most frequent
(``top'') as well as how many times it appears (``freq''):

\begin{itemize}
\tightlist
\item
  Philadelphia and Chicago remain primarily ``Low'' price bins,
  suggesting generally lower median prices compared to the other metros.
\item
  Boston and DC hover around the ``High'' bin most often, indicating a
  consistently elevated market.
\item
  LA and Seattle both stand out with ``Very High'' as their predominant
  bin, showcasing top-tier market prices.
\item
  National data is skewed to ``Low'', which makes sense as a broad
  average across diverse regions---including many less expensive areas
  beyond major metro hubs.
\end{itemize}

Overall, these frequency counts confirm each metro's price tier in the
dataset, supporting the broader trends observed in the scatterplots and
boxplots (e.g., LA's and Seattle's especially high prices,
Philadelphia's lower cost, etc.).

\begin{center}\rule{0.5\linewidth}{0.5pt}\end{center}

    \begin{tcolorbox}[breakable, size=fbox, boxrule=1pt, pad at break*=1mm,colback=cellbackground, colframe=cellborder]
\prompt{In}{incolor}{14}{\boxspacing}
\begin{Verbatim}[commandchars=\\\{\}]
\PY{n}{e}\PY{o}{.}\PY{n}{plot\PYZus{}average\PYZus{}price\PYZus{}over\PYZus{}time}\PY{p}{(}\PY{n}{redfin\PYZus{}df}\PY{p}{)}
\end{Verbatim}
\end{tcolorbox}

    \begin{center}
    \adjustimage{max size={0.9\linewidth}{0.9\paperheight}}{thesis_project_files/thesis_project_29_0.png}
    \end{center}
    { \hspace*{\fill} \\}
    
    Average Price Over Time Insights

Below is an interpretation of the ``Average Price Over Time'' line chart
and how it confirms or extends what we've seen in the other plots.

\begin{enumerate}
\def\labelenumi{\arabic{enumi}.}
\tightlist
\item
  Overall Upward Trajectory\\
  From about \textbackslash\$250k in 2012 to roughly \$550k+ by late
  2024, the chart depicts a steady, long-term increase in median sale
  prices at the national level.

  \begin{itemize}
  \item
    This upward climb aligns with the Nationwide ``Low'' bin we saw in
    price\_bin distributions; although it's labeled ``Low'' relative to
    high-cost metros like Los Angeles or Seattle, the absolute trend is
    still consistently rising.
  \item ~
    \subsection{It also matches our time-series decomposition results
    (and scatterplots) indicating that price dips were mostly
    short-lived despite fluctuations in inventory or days on
    market.}\label{it-also-matches-our-time-series-decomposition-results-and-scatterplots-indicating-that-price-dips-were-mostly-short-lived-despite-fluctuations-in-inventory-or-days-on-market.}
  \end{itemize}
\item
  Mid-2019 to 2021 Acceleration Around 2019, the slope appears to
  steepen briefly, followed by a small plateau in 2023--2024.

  \begin{itemize}
  \item
    Some earlier scatterplots (e.g., inventory\_mom
    vs.~median\_sale\_price) showed that surges in inventory didn't
    immediately reduce prices---this line chart confirms a continued
    upward price movement during that period.
  \item ~
    \subsection{The slight plateau in 2023--2024 could tie back to the
    occasional inventory spikes or days-on-market increases we saw for
    certain metros, suggesting a softening demand or a leveling out in
    the national
    average.}\label{the-slight-plateau-in-20232024-could-tie-back-to-the-occasional-inventory-spikes-or-days-on-market-increases-we-saw-for-certain-metros-suggesting-a-softening-demand-or-a-leveling-out-in-the-national-average.}
  \end{itemize}
\item
  Tying Back to Region-Level Diversity\\
  While this chart aggregates the National data, we know from the
  region-based boxplots and scatterplots that individual metros like
  Boston or Los Angeles often exceed these values, whereas places like
  Chicago or Philadelphia consistently fall below.

  \begin{itemize}
  \item ~
    \subsection{The overall national upward line masks the substantial
    variations visible in local scatterplots. L.A. (Red) or Seattle
    (Brown) can be ``Very High,'' while Philly (Purple) is mostly
    ``Low,'' yet everyone is trending
    upward.}\label{the-overall-national-upward-line-masks-the-substantial-variations-visible-in-local-scatterplots.-l.a.-red-or-seattle-brown-can-be-very-high-while-philly-purple-is-mostly-low-yet-everyone-is-trending-upward.}

    Final Takeaway\\
    This time-series line confirms a broad, long-run price appreciation,
    reinforcing that local anomalies (whether spikes in inventory or
    changes in days on market) have not derailed the national climb. It
    fits well with the earlier findings that supply surges sometimes lag
    before affecting prices and that regional differences are key, even
    within an overall upward trend.
  \end{itemize}
\end{enumerate}

\begin{center}\rule{0.5\linewidth}{0.5pt}\end{center}

    \subsubsection{Week 3 EDA}\label{week-3-eda}

Week 3: - Tasks: - Perform feature engineering. - Handle missing values
and outliers. - Goal: Prepare a well-structured dataset for modeling.

    \paragraph{Outliers}\label{outliers}

    Why Are We Checking for Outliers?

The goal of this step is to \textbf{identify and handle extreme values}
that could distort our analysis and model performance. Outliers can
arise from \textbf{data entry errors, rare events, or genuine market
fluctuations}, and it's crucial to assess their impact.

\begin{itemize}
\tightlist
\item
  \textbf{Why It Matters:}

  \begin{itemize}
  \tightlist
  \item
    Outliers can \textbf{skew statistical summaries} (e.g., mean,
    standard deviation).\\
  \item
    They may \textbf{negatively impact model training}, especially in
    regression and machine learning algorithms.\\
  \item
    Identifying them helps decide whether to \textbf{remove, transform,
    or cap} extreme values for better data integrity.
  \end{itemize}
\end{itemize}

By detecting outliers, we ensure that our dataset remains
\textbf{accurate, reliable, and representative} of real-world trends.

    \begin{tcolorbox}[breakable, size=fbox, boxrule=1pt, pad at break*=1mm,colback=cellbackground, colframe=cellborder]
\prompt{In}{incolor}{15}{\boxspacing}
\begin{Verbatim}[commandchars=\\\{\}]
\PY{c+c1}{\PYZsh{} Checking for Outliers}
\PY{n}{outliers} \PY{o}{=} \PY{n}{e}\PY{o}{.}\PY{n}{detect\PYZus{}outliers\PYZus{}iqr}\PY{p}{(}\PY{n}{redfin\PYZus{}df}\PY{p}{)}
\PY{n+nb}{print}\PY{p}{(}\PY{n}{outliers}\PY{p}{)}
\end{Verbatim}
\end{tcolorbox}

            \begin{tcolorbox}[breakable, size=fbox, boxrule=.5pt, pad at break*=1mm, opacityfill=0]
\prompt{Out}{outcolor}{15}{\boxspacing}
\begin{Verbatim}[commandchars=\\\{\}]
median\_sale\_price             0
median\_sale\_price\_mom        27
median\_sale\_price\_yoy        37
homes\_sold                  157
homes\_sold\_mom               17
homes\_sold\_yoy               62
new\_listings                157
new\_listings\_mom             79
new\_listings\_yoy             88
inventory                   157
inventory\_mom                84
inventory\_yoy                46
days\_on\_market               29
days\_on\_market\_mom           80
days\_on\_market\_yoy           75
average\_sale\_to\_list         41
average\_sale\_to\_list\_mom     50
average\_sale\_to\_list\_yoy     97
dtype: int64
\end{Verbatim}
\end{tcolorbox}
        
    Outlier Insights

Several columns have notable outliers, particularly in: - homes\_sold
(157 outliers), new\_listings (157 outliers), and inventory (157
outliers): These metrics likely have extreme values in certain months or
regions. - median\_sale\_price\_yoy (37 outliers) and homes\_sold\_yoy
(62 outliers): High year-over-year changes indicate significant
fluctuations in the market. - days\_on\_market\_mom (80 outliers) and
days\_on\_market\_yoy (75 outliers): The time a home stays on the market
shows variability. - average\_sale\_to\_list\_yoy (97 outliers): Changes
in sale-to-list price ratios suggest market volatility.

\begin{center}\rule{0.5\linewidth}{0.5pt}\end{center}

    \begin{tcolorbox}[breakable, size=fbox, boxrule=1pt, pad at break*=1mm,colback=cellbackground, colframe=cellborder]
\prompt{In}{incolor}{ }{\boxspacing}
\begin{Verbatim}[commandchars=\\\{\}]
\PY{c+c1}{\PYZsh{} Selected columns with significant outliers}
\PY{n}{outlier\PYZus{}cols} \PY{o}{=} \PY{p}{[}\PY{l+s+s2}{\PYZdq{}}\PY{l+s+s2}{homes\PYZus{}sold}\PY{l+s+s2}{\PYZdq{}}\PY{p}{,} 
                \PY{l+s+s2}{\PYZdq{}}\PY{l+s+s2}{new\PYZus{}listings}\PY{l+s+s2}{\PYZdq{}}\PY{p}{,} 
                \PY{l+s+s2}{\PYZdq{}}\PY{l+s+s2}{inventory}\PY{l+s+s2}{\PYZdq{}}\PY{p}{,} 
                \PY{l+s+s2}{\PYZdq{}}\PY{l+s+s2}{median\PYZus{}sale\PYZus{}price\PYZus{}yoy}\PY{l+s+s2}{\PYZdq{}}\PY{p}{,} 
                \PY{l+s+s2}{\PYZdq{}}\PY{l+s+s2}{homes\PYZus{}sold\PYZus{}yoy}\PY{l+s+s2}{\PYZdq{}}\PY{p}{,} 
                \PY{l+s+s2}{\PYZdq{}}\PY{l+s+s2}{days\PYZus{}on\PYZus{}market\PYZus{}mom}\PY{l+s+s2}{\PYZdq{}}\PY{p}{,} 
                \PY{l+s+s2}{\PYZdq{}}\PY{l+s+s2}{days\PYZus{}on\PYZus{}market\PYZus{}yoy}\PY{l+s+s2}{\PYZdq{}}\PY{p}{,} 
                \PY{l+s+s2}{\PYZdq{}}\PY{l+s+s2}{average\PYZus{}sale\PYZus{}to\PYZus{}list\PYZus{}yoy}\PY{l+s+s2}{\PYZdq{}}\PY{p}{]}
\end{Verbatim}
\end{tcolorbox}

    \begin{tcolorbox}[breakable, size=fbox, boxrule=1pt, pad at break*=1mm,colback=cellbackground, colframe=cellborder]
\prompt{In}{incolor}{16}{\boxspacing}
\begin{Verbatim}[commandchars=\\\{\}]
\PY{c+c1}{\PYZsh{} Visualizing Outliers}
\PY{n}{e}\PY{o}{.}\PY{n}{visualize\PYZus{}outliers}\PY{p}{(}\PY{n}{redfin\PYZus{}df}\PY{p}{,} \PY{n}{outlier\PYZus{}cols}\PY{p}{)}
\end{Verbatim}
\end{tcolorbox}

    \begin{center}
    \adjustimage{max size={0.9\linewidth}{0.9\paperheight}}{thesis_project_files/thesis_project_37_0.png}
    \end{center}
    { \hspace*{\fill} \\}
    
    Boxplot of Columns with Significant Outliers Insights

The boxplot confirms the presence of extreme outliers, particularly in
homes\_sold, new\_listings, and inventory. These could be due to market
fluctuations or anomalies in reporting.

\begin{center}\rule{0.5\linewidth}{0.5pt}\end{center}

    Handling Strategies:

\begin{enumerate}
\def\labelenumi{\arabic{enumi}.}
\tightlist
\item
  Winsorization: Capping extreme values at the 5th and 95th percentiles
  to reduce impact.\\
\item
  Log Transformation: Applying a log transformation to reduce skewness
  (useful if data is highly right-skewed).\\
\item
  Outlier Removal: If values seem unrealistic or errors, we can remove
  them.
\end{enumerate}

I'll apply Winsorization (capping extreme values) for better
distribution while preserving trends.

\begin{center}\rule{0.5\linewidth}{0.5pt}\end{center}

    \begin{tcolorbox}[breakable, size=fbox, boxrule=1pt, pad at break*=1mm,colback=cellbackground, colframe=cellborder]
\prompt{In}{incolor}{36}{\boxspacing}
\begin{Verbatim}[commandchars=\\\{\}]
\PY{n}{df\PYZus{}winsorized} \PY{o}{=} \PY{n}{e}\PY{o}{.}\PY{n}{apply\PYZus{}winsorization}\PY{p}{(}\PY{n}{redfin\PYZus{}df}\PY{p}{,} \PY{n}{outlier\PYZus{}cols}\PY{p}{)}
\end{Verbatim}
\end{tcolorbox}

    \begin{center}
    \adjustimage{max size={0.9\linewidth}{0.9\paperheight}}{thesis_project_files/thesis_project_40_0.png}
    \end{center}
    { \hspace*{\fill} \\}
    
    Boxplot After Winsorization Insight

After Winsorization, the extreme outliers have been capped, making the
distribution more balanced while retaining market trends.

\begin{center}\rule{0.5\linewidth}{0.5pt}\end{center}

    Why Are We Creating a Correlation Matrix Again?

The goal of this step is to \textbf{reassess the relationships between
features after Winsorization}, ensuring that our data transformations
have not significantly altered key correlations.

\begin{itemize}
\tightlist
\item
  \textbf{Why It Matters:}

  \begin{itemize}
  \tightlist
  \item
    Winsorization \textbf{caps extreme values} to reduce the impact of
    outliers. This can subtly shift correlations between variables.\\
  \item
    We need to confirm whether the \textbf{same strong relationships
    still exist} after handling outliers.\\
  \item
    If correlations remain unchanged, it validates that
    \textbf{Winsorization preserved overall trends} while mitigating the
    effects of extreme values.\\
  \item
    If correlations change significantly, we may need to revisit
    \textbf{feature selection} and \textbf{data transformations} to
    ensure model reliability.
  \end{itemize}
\end{itemize}

By computing the correlation matrix post-Winsorization, we verify that
our dataset remains \textbf{structurally sound for modeling and
analysis.}

    \begin{tcolorbox}[breakable, size=fbox, boxrule=1pt, pad at break*=1mm,colback=cellbackground, colframe=cellborder]
\prompt{In}{incolor}{22}{\boxspacing}
\begin{Verbatim}[commandchars=\\\{\}]
\PY{c+c1}{\PYZsh{} Compute correlation matrix}
\PY{n}{correlation\PYZus{}matrix} \PY{o}{=} \PY{n}{df\PYZus{}winsorized}\PY{o}{.}\PY{n}{corr}\PY{p}{(}\PY{p}{)}
\end{Verbatim}
\end{tcolorbox}

    \begin{tcolorbox}[breakable, size=fbox, boxrule=1pt, pad at break*=1mm,colback=cellbackground, colframe=cellborder]
\prompt{In}{incolor}{23}{\boxspacing}
\begin{Verbatim}[commandchars=\\\{\}]
\PY{k+kn}{from} \PY{n+nn}{IPython}\PY{n+nn}{.}\PY{n+nn}{display} \PY{k+kn}{import} \PY{n}{display}
\PY{n}{display}\PY{p}{(}\PY{n}{correlation\PYZus{}matrix}\PY{p}{)}
\end{Verbatim}
\end{tcolorbox}

    
    \begin{Verbatim}[commandchars=\\\{\}]
                          median\_sale\_price  median\_sale\_price\_mom  \textbackslash{}
median\_sale\_price                  1.000000               0.005228   
median\_sale\_price\_mom              0.005228               1.000000   
median\_sale\_price\_yoy              0.098092               0.178894   
homes\_sold                        -0.232658               0.014147   
homes\_sold\_mom                     0.001816               0.653534   
homes\_sold\_yoy                    -0.265236               0.077921   
new\_listings                      -0.237221               0.018802   
new\_listings\_mom                  -0.012778               0.015628   
new\_listings\_yoy                  -0.076005               0.053991   
inventory                         -0.270464              -0.002811   
inventory\_mom                      0.079763               0.205451   
inventory\_yoy                      0.171810              -0.101212   
days\_on\_market                    -0.472584              -0.075931   
days\_on\_market\_mom                 0.036927              -0.626741   
days\_on\_market\_yoy                 0.236736              -0.103519   
average\_sale\_to\_list               0.758161               0.111648   
average\_sale\_to\_list\_mom          -0.055247               0.556545   
average\_sale\_to\_list\_yoy          -0.067936               0.129414   

                          median\_sale\_price\_yoy  homes\_sold  homes\_sold\_mom  \textbackslash{}
median\_sale\_price                      0.098092   -0.232658        0.001816   
median\_sale\_price\_mom                  0.178894    0.014147        0.653534   
median\_sale\_price\_yoy                  1.000000    0.057881        0.022894   
homes\_sold                             0.057881    1.000000        0.014527   
homes\_sold\_mom                         0.022894    0.014527        1.000000   
homes\_sold\_yoy                         0.324370    0.024509        0.127307   
new\_listings                           0.052116    0.985035        0.018705   
new\_listings\_mom                      -0.018807   -0.041356       -0.024779   
new\_listings\_yoy                       0.295929    0.018205        0.114331   
inventory                              0.020582    0.954897       -0.004080   
inventory\_mom                          0.057880    0.000578        0.234952   
inventory\_yoy                         -0.285196    0.000890       -0.047563   
days\_on\_market                        -0.276977    0.120669       -0.090630   
days\_on\_market\_mom                    -0.069687   -0.032259       -0.608613   
days\_on\_market\_yoy                    -0.418244   -0.048649       -0.031243   
average\_sale\_to\_list                   0.387972   -0.084859        0.113571   
average\_sale\_to\_list\_mom               0.090790   -0.003010        0.496990   
average\_sale\_to\_list\_yoy               0.610831   -0.019541        0.052232   

                          homes\_sold\_yoy  new\_listings  new\_listings\_mom  \textbackslash{}
median\_sale\_price              -0.265236     -0.237221         -0.012778   
median\_sale\_price\_mom           0.077921      0.018802          0.015628   
median\_sale\_price\_yoy           0.324370      0.052116         -0.018807   
homes\_sold                      0.024509      0.985035         -0.041356   
homes\_sold\_mom                  0.127307      0.018705         -0.024779   
homes\_sold\_yoy                  1.000000      0.024207          0.008135   
new\_listings                    0.024207      1.000000         -0.004320   
new\_listings\_mom                0.008135     -0.004320          1.000000   
new\_listings\_yoy                0.587626      0.022839          0.091386   
inventory                       0.039470      0.958168         -0.025606   
inventory\_mom                  -0.069369      0.032035          0.482830   
inventory\_yoy                  -0.336885      0.002229         -0.013579   
days\_on\_market                  0.130505      0.134936          0.180043   
days\_on\_market\_mom             -0.079609     -0.052766         -0.061885   
days\_on\_market\_yoy             -0.532169     -0.050926          0.004886   
average\_sale\_to\_list           -0.161861     -0.090729         -0.035927   
average\_sale\_to\_list\_mom        0.170576      0.010565          0.215508   
average\_sale\_to\_list\_yoy        0.567886     -0.019728          0.003995   

                          new\_listings\_yoy  inventory  inventory\_mom  \textbackslash{}
median\_sale\_price                -0.076005  -0.270464       0.079763   
median\_sale\_price\_mom             0.053991  -0.002811       0.205451   
median\_sale\_price\_yoy             0.295929   0.020582       0.057880   
homes\_sold                        0.018205   0.954897       0.000578   
homes\_sold\_mom                    0.114331  -0.004080       0.234952   
homes\_sold\_yoy                    0.587626   0.039470      -0.069369   
new\_listings                      0.022839   0.958168       0.032035   
new\_listings\_mom                  0.091386  -0.025606       0.482830   
new\_listings\_yoy                  1.000000   0.021441       0.144859   
inventory                         0.021441   1.000000      -0.002156   
inventory\_mom                     0.144859  -0.002156       1.000000   
inventory\_yoy                    -0.085640   0.013893       0.087980   
days\_on\_market                   -0.058362   0.207376      -0.148243   
days\_on\_market\_mom               -0.044884  -0.019536      -0.328943   
days\_on\_market\_yoy               -0.303271  -0.054352       0.006764   
average\_sale\_to\_list              0.026335  -0.155334       0.243446   
average\_sale\_to\_list\_mom          0.131237  -0.012001       0.200081   
average\_sale\_to\_list\_yoy          0.397182  -0.038531       0.029244   

                          inventory\_yoy  days\_on\_market  days\_on\_market\_mom  \textbackslash{}
median\_sale\_price              0.171810       -0.472584            0.036927   
median\_sale\_price\_mom         -0.101212       -0.075931           -0.626741   
median\_sale\_price\_yoy         -0.285196       -0.276977           -0.069687   
homes\_sold                     0.000890        0.120669           -0.032259   
homes\_sold\_mom                -0.047563       -0.090630           -0.608613   
homes\_sold\_yoy                -0.336885        0.130505           -0.079609   
new\_listings                   0.002229        0.134936           -0.052766   
new\_listings\_mom              -0.013579        0.180043           -0.061885   
new\_listings\_yoy              -0.085640       -0.058362           -0.044884   
inventory                      0.013893        0.207376           -0.019536   
inventory\_mom                  0.087980       -0.148243           -0.328943   
inventory\_yoy                  1.000000        0.012950            0.131862   
days\_on\_market                 0.012950        1.000000            0.114098   
days\_on\_market\_mom             0.131862        0.114098            1.000000   
days\_on\_market\_yoy             0.585719       -0.052597            0.196396   
average\_sale\_to\_list          -0.067417       -0.684080           -0.078461   
average\_sale\_to\_list\_mom      -0.211184        0.014321           -0.551240   
average\_sale\_to\_list\_yoy      -0.611348       -0.082182           -0.118623   

                          days\_on\_market\_yoy  average\_sale\_to\_list  \textbackslash{}
median\_sale\_price                   0.236736              0.758161   
median\_sale\_price\_mom              -0.103519              0.111648   
median\_sale\_price\_yoy              -0.418244              0.387972   
homes\_sold                         -0.048649             -0.084859   
homes\_sold\_mom                     -0.031243              0.113571   
homes\_sold\_yoy                     -0.532169             -0.161861   
new\_listings                       -0.050926             -0.090729   
new\_listings\_mom                    0.004886             -0.035927   
new\_listings\_yoy                   -0.303271              0.026335   
inventory                          -0.054352             -0.155334   
inventory\_mom                       0.006764              0.243446   
inventory\_yoy                       0.585719             -0.067417   
days\_on\_market                     -0.052597             -0.684080   
days\_on\_market\_mom                  0.196396             -0.078461   
days\_on\_market\_yoy                  1.000000              0.078648   
average\_sale\_to\_list                0.078648              1.000000   
average\_sale\_to\_list\_mom           -0.141950              0.080975   
average\_sale\_to\_list\_yoy           -0.695177              0.214099   

                          average\_sale\_to\_list\_mom  average\_sale\_to\_list\_yoy  
median\_sale\_price                        -0.055247                 -0.067936  
median\_sale\_price\_mom                     0.556545                  0.129414  
median\_sale\_price\_yoy                     0.090790                  0.610831  
homes\_sold                               -0.003010                 -0.019541  
homes\_sold\_mom                            0.496990                  0.052232  
homes\_sold\_yoy                            0.170576                  0.567886  
new\_listings                              0.010565                 -0.019728  
new\_listings\_mom                          0.215508                  0.003995  
new\_listings\_yoy                          0.131237                  0.397182  
inventory                                -0.012001                 -0.038531  
inventory\_mom                             0.200081                  0.029244  
inventory\_yoy                            -0.211184                 -0.611348  
days\_on\_market                            0.014321                 -0.082182  
days\_on\_market\_mom                       -0.551240                 -0.118623  
days\_on\_market\_yoy                       -0.141950                 -0.695177  
average\_sale\_to\_list                      0.080975                  0.214099  
average\_sale\_to\_list\_mom                  1.000000                  0.261101  
average\_sale\_to\_list\_yoy                  0.261101                  1.000000  
    \end{Verbatim}

    
    \begin{tcolorbox}[breakable, size=fbox, boxrule=1pt, pad at break*=1mm,colback=cellbackground, colframe=cellborder]
\prompt{In}{incolor}{30}{\boxspacing}
\begin{Verbatim}[commandchars=\\\{\}]
\PY{n}{e}\PY{o}{.}\PY{n}{plot\PYZus{}correlation\PYZus{}heatmap}\PY{p}{(}\PY{n}{redfin\PYZus{}df}\PY{p}{)}
\end{Verbatim}
\end{tcolorbox}

    \begin{center}
    \adjustimage{max size={0.9\linewidth}{0.9\paperheight}}{thesis_project_files/thesis_project_45_0.png}
    \end{center}
    { \hspace*{\fill} \\}
    
    Key Findings from Correlation Analysis

Since Winsorization did not significantly impact the correlation
structure, it suggests that the extreme values, while present, were not
highly influential in distorting the relationships between variables.
This means the dataset's original trends were robust even with the
outliers.

\begin{itemize}
\tightlist
\item
  The high correlation between homes\_sold and new\_listings (0.985)
  suggests that one of these could be dropped without significant
  information loss.
\item
  inventory remains strongly correlated with homes\_sold (0.955),
  implying potential redundancy.
\item
  Market trends still hold:

  \begin{itemize}
  \tightlist
  \item
    Higher prices → Faster sales (Negative correlation between
    median\_sale\_price \& days\_on\_market: -0.473).
  \item
    Price increases reduce days on market (Negative correlation between
    median\_sale\_price\_yoy \& days\_on\_market\_yoy: -0.418).
  \item
    Sale-to-list ratio trends with median prices (0.611 correlation
    between average\_sale\_to\_list\_yoy \& median\_sale\_price\_yoy).
  \end{itemize}
\end{itemize}

\begin{center}\rule{0.5\linewidth}{0.5pt}\end{center}

    \paragraph{Feature Engineering}\label{feature-engineering}

    Feature Engineering Discussion

The goal of this step is to create meaningful features that enhance the
dataset's ability to capture market trends and improve model
performance. By engineering new features, we aim to make patterns more
interpretable and informative for downstream analysis.

\begin{itemize}
\tightlist
\item
  Market Demand Ratio (homes\_sold / new\_listings)

  \begin{itemize}
  \tightlist
  \item
    Helps measure market competitiveness by indicating how quickly homes
    are being sold relative to new listings.
  \item
    A higher ratio suggests a strong seller's market, while a lower
    ratio may indicate higher inventory relative to demand.
  \end{itemize}
\item
  Days on Market Categorization (Fast, Moderate, Slow)

  \begin{itemize}
  \tightlist
  \item
    Transforms days\_on\_market from a raw numeric value into categories
    that are easier to analyze.
  \item
    This allows us to better understand how quickly homes are selling
    and whether certain market conditions lead to longer selling times.
  \end{itemize}
\end{itemize}

By performing feature engineering at this stage, we are ensuring that
the dataset is not just clean, but also structured in a way that makes
it more valuable for analysis and predictive modeling.

\begin{center}\rule{0.5\linewidth}{0.5pt}\end{center}

    \begin{tcolorbox}[breakable, size=fbox, boxrule=1pt, pad at break*=1mm,colback=cellbackground, colframe=cellborder]
\prompt{In}{incolor}{33}{\boxspacing}
\begin{Verbatim}[commandchars=\\\{\}]
\PY{c+c1}{\PYZsh{} Feature Engineering}
\PY{n}{df\PYZus{}transformed} \PY{o}{=} \PY{n}{e}\PY{o}{.}\PY{n}{create\PYZus{}features}\PY{p}{(}\PY{n}{df\PYZus{}winsorized}\PY{p}{)}
\PY{n}{df\PYZus{}transformed}\PY{o}{.}\PY{n}{head}\PY{p}{(}\PY{p}{)}
\end{Verbatim}
\end{tcolorbox}

            \begin{tcolorbox}[breakable, size=fbox, boxrule=.5pt, pad at break*=1mm, opacityfill=0]
\prompt{Out}{outcolor}{33}{\boxspacing}
\begin{Verbatim}[commandchars=\\\{\}]
      region month\_of\_period\_end  median\_sale\_price  median\_sale\_price\_mom  \textbackslash{}
0   National          2022-07-01             413000                   -3.6
1   National          2022-08-01             407000                   -1.5
2   National          2022-06-01             429000                   -0.7
3   National          2021-08-01             381000                   -1.1
4   National          2021-09-01             377000                   -1.0

   median\_sale\_price\_yoy  homes\_sold  homes\_sold\_mom  homes\_sold\_yoy  \textbackslash{}
0                    7.4      530476           -14.4           -21.9
1                    6.9      530476             5.4           -16.9
2                   10.8      530476             4.1           -15.0
3                   16.2      530476            -1.0             0.1
4                   13.9      530476            -5.1            -4.7

   new\_listings  new\_listings\_mom  {\ldots}  inventory\_yoy  days\_on\_market  \textbackslash{}
0        656191             -14.3  {\ldots}           27.1              21
1        622537             -11.7  {\ldots}           26.7              27
2        656191               6.4  {\ldots}           28.1              18
3        656191              -9.8  {\ldots}          -18.2              17
4        656191              -6.5  {\ldots}          -16.1              19

   days\_on\_market\_mom  days\_on\_market\_yoy  average\_sale\_to\_list  \textbackslash{}
0                   3                   6                 101.0
1                   5                  10                  99.9
2                   1                   3                 102.2
3                   2                 -15                 101.6
4                   2                 -11                 101.0

   average\_sale\_to\_list\_mom  average\_sale\_to\_list\_yoy  price\_bin  \textbackslash{}
0                      -1.2                      -1.2       High
1                      -1.1                      -1.7       High
2                      -0.8                      -0.3       High
3                      -0.6                       2.1       High
4                      -0.6                       1.7       High

   market\_demand\_ratio  days\_on\_market\_category
0             0.808417                 Moderate
1             0.852120                 Moderate
2             0.808417                 Moderate
3             0.808417                 Moderate
4             0.808417                 Moderate

[5 rows x 23 columns]
\end{Verbatim}
\end{tcolorbox}
        
    \paragraph{Feature Selection}\label{feature-selection}

    Feature Selection and Reduction Discussion

The goal here is to reduce redundancy in the dataset by removing highly
correlated features while retaining the most informative ones for
modeling.

\begin{itemize}
\tightlist
\item
  Identifying Redundant Features From our correlation matrix:

  \begin{itemize}
  \tightlist
  \item
    homes\_sold vs.~new\_listings (0.985 correlation)

    \begin{itemize}
    \tightlist
    \item
      Implication: Almost a direct relationship---regions with higher
      home sales tend to have higher new listings.
    \item
      Decision: Drop one---likely new\_listings because homes\_sold
      directly reflects actual transactions rather than just
      availability.
    \end{itemize}
  \item
    inventory vs.~homes\_sold (0.955 correlation)

    \begin{itemize}
    \tightlist
    \item
      Implication: Inventory reflects available homes, but homes sold
      captures actual transactions.
    \item
      Decision: Drop one---likely inventory since homes\_sold is more
      action-based.
    \end{itemize}
  \end{itemize}
\end{itemize}

\begin{center}\rule{0.5\linewidth}{0.5pt}\end{center}

    \begin{tcolorbox}[breakable, size=fbox, boxrule=1pt, pad at break*=1mm,colback=cellbackground, colframe=cellborder]
\prompt{In}{incolor}{35}{\boxspacing}
\begin{Verbatim}[commandchars=\\\{\}]
\PY{c+c1}{\PYZsh{} Drop redundant features: \PYZsq{}new\PYZus{}listings\PYZsq{} and \PYZsq{}inventory\PYZsq{}}
\PY{n}{df\PYZus{}reduced} \PY{o}{=} \PY{n}{e}\PY{o}{.}\PY{n}{drop\PYZus{}redundant\PYZus{}features}\PY{p}{(}\PY{n}{df\PYZus{}winsorized}\PY{p}{,} \PY{p}{[}\PY{l+s+s1}{\PYZsq{}}\PY{l+s+s1}{new\PYZus{}listings}\PY{l+s+s1}{\PYZsq{}}\PY{p}{,} \PY{l+s+s1}{\PYZsq{}}\PY{l+s+s1}{inventory}\PY{l+s+s1}{\PYZsq{}}\PY{p}{]}\PY{p}{)}

\PY{c+c1}{\PYZsh{} Display the updated dataset after feature reduction}
\PY{n}{df\PYZus{}reduced}\PY{o}{.}\PY{n}{head}\PY{p}{(}\PY{p}{)}
\end{Verbatim}
\end{tcolorbox}

            \begin{tcolorbox}[breakable, size=fbox, boxrule=.5pt, pad at break*=1mm, opacityfill=0]
\prompt{Out}{outcolor}{35}{\boxspacing}
\begin{Verbatim}[commandchars=\\\{\}]
      region month\_of\_period\_end  median\_sale\_price  median\_sale\_price\_mom  \textbackslash{}
0   National          2022-07-01             413000                   -3.6
1   National          2022-08-01             407000                   -1.5
2   National          2022-06-01             429000                   -0.7
3   National          2021-08-01             381000                   -1.1
4   National          2021-09-01             377000                   -1.0

   median\_sale\_price\_yoy  homes\_sold  homes\_sold\_mom  homes\_sold\_yoy  \textbackslash{}
0                    7.4      530476           -14.4           -21.9
1                    6.9      530476             5.4           -16.9
2                   10.8      530476             4.1           -15.0
3                   16.2      530476            -1.0             0.1
4                   13.9      530476            -5.1            -4.7

   new\_listings\_mom  new\_listings\_yoy  {\ldots}  inventory\_yoy  days\_on\_market  \textbackslash{}
0             -14.3             -11.9  {\ldots}           27.1              21
1             -11.7             -13.7  {\ldots}           26.7              27
2               6.4               1.0  {\ldots}           28.1              18
3              -9.8               0.0  {\ldots}          -18.2              17
4              -6.5              -4.5  {\ldots}          -16.1              19

   days\_on\_market\_mom  days\_on\_market\_yoy  average\_sale\_to\_list  \textbackslash{}
0                   3                   6                 101.0
1                   5                  10                  99.9
2                   1                   3                 102.2
3                   2                 -15                 101.6
4                   2                 -11                 101.0

   average\_sale\_to\_list\_mom  average\_sale\_to\_list\_yoy  price\_bin  \textbackslash{}
0                      -1.2                      -1.2       High
1                      -1.1                      -1.7       High
2                      -0.8                      -0.3       High
3                      -0.6                       2.1       High
4                      -0.6                       1.7       High

  market\_demand\_ratio  days\_on\_market\_category
0            0.808417                 Moderate
1            0.852120                 Moderate
2            0.808417                 Moderate
3            0.808417                 Moderate
4            0.808417                 Moderate

[5 rows x 21 columns]
\end{Verbatim}
\end{tcolorbox}
        
    \subsubsection{Week 4:}\label{week-4}

\begin{itemize}
\tightlist
\item
  Tasks:

  \begin{itemize}
  \tightlist
  \item
    Select baseline models for initial testing.
  \item
    Train and evaluate simple models.
  \end{itemize}
\item
  Goal: Identify the best initial modeling approach.
\end{itemize}

    \begin{tcolorbox}[breakable, size=fbox, boxrule=1pt, pad at break*=1mm,colback=cellbackground, colframe=cellborder]
\prompt{In}{incolor}{ }{\boxspacing}
\begin{Verbatim}[commandchars=\\\{\}]

\end{Verbatim}
\end{tcolorbox}

    \begin{tcolorbox}[breakable, size=fbox, boxrule=1pt, pad at break*=1mm,colback=cellbackground, colframe=cellborder]
\prompt{In}{incolor}{ }{\boxspacing}
\begin{Verbatim}[commandchars=\\\{\}]

\end{Verbatim}
\end{tcolorbox}

    \begin{tcolorbox}[breakable, size=fbox, boxrule=1pt, pad at break*=1mm,colback=cellbackground, colframe=cellborder]
\prompt{In}{incolor}{ }{\boxspacing}
\begin{Verbatim}[commandchars=\\\{\}]

\end{Verbatim}
\end{tcolorbox}

    \begin{tcolorbox}[breakable, size=fbox, boxrule=1pt, pad at break*=1mm,colback=cellbackground, colframe=cellborder]
\prompt{In}{incolor}{ }{\boxspacing}
\begin{Verbatim}[commandchars=\\\{\}]

\end{Verbatim}
\end{tcolorbox}

    \begin{tcolorbox}[breakable, size=fbox, boxrule=1pt, pad at break*=1mm,colback=cellbackground, colframe=cellborder]
\prompt{In}{incolor}{ }{\boxspacing}
\begin{Verbatim}[commandchars=\\\{\}]

\end{Verbatim}
\end{tcolorbox}

    \begin{tcolorbox}[breakable, size=fbox, boxrule=1pt, pad at break*=1mm,colback=cellbackground, colframe=cellborder]
\prompt{In}{incolor}{ }{\boxspacing}
\begin{Verbatim}[commandchars=\\\{\}]

\end{Verbatim}
\end{tcolorbox}

    \begin{tcolorbox}[breakable, size=fbox, boxrule=1pt, pad at break*=1mm,colback=cellbackground, colframe=cellborder]
\prompt{In}{incolor}{ }{\boxspacing}
\begin{Verbatim}[commandchars=\\\{\}]

\end{Verbatim}
\end{tcolorbox}

    \begin{tcolorbox}[breakable, size=fbox, boxrule=1pt, pad at break*=1mm,colback=cellbackground, colframe=cellborder]
\prompt{In}{incolor}{ }{\boxspacing}
\begin{Verbatim}[commandchars=\\\{\}]

\end{Verbatim}
\end{tcolorbox}

    \begin{tcolorbox}[breakable, size=fbox, boxrule=1pt, pad at break*=1mm,colback=cellbackground, colframe=cellborder]
\prompt{In}{incolor}{ }{\boxspacing}
\begin{Verbatim}[commandchars=\\\{\}]

\end{Verbatim}
\end{tcolorbox}

    \begin{tcolorbox}[breakable, size=fbox, boxrule=1pt, pad at break*=1mm,colback=cellbackground, colframe=cellborder]
\prompt{In}{incolor}{ }{\boxspacing}
\begin{Verbatim}[commandchars=\\\{\}]

\end{Verbatim}
\end{tcolorbox}

    \begin{tcolorbox}[breakable, size=fbox, boxrule=1pt, pad at break*=1mm,colback=cellbackground, colframe=cellborder]
\prompt{In}{incolor}{ }{\boxspacing}
\begin{Verbatim}[commandchars=\\\{\}]

\end{Verbatim}
\end{tcolorbox}


    % Add a bibliography block to the postdoc
    
    
    
\end{document}
